% !TeX program = pdflatex
% Resumen ejecutivo de la propuesta curricular (formato OPES--CONARE).
% Compilar desde esta carpeta: python scripts/generate_resumen_ejecutivo.py (regenera las tablas de los apartados 8 a 10)
% y luego pdflatex resumen-ejecutivo (dos veces).
\documentclass[11pt,a4paper]{article}

\usepackage[utf8]{inputenc}
\usepackage[T1]{fontenc}
\usepackage{lmodern}
\renewcommand{\familydefault}{\sfdefault}
\usepackage[spanish, es-tabla]{babel}
\usepackage[left=1.9cm,right=1.9cm,top=2.2cm,bottom=2.2cm,headheight=14pt]{geometry}
\usepackage[svgnames, dvipsnames, table]{xcolor}
\usepackage{amsmath, amsfonts, amssymb}
\usepackage{graphicx}
\usepackage{tikz}
\usetikzlibrary{calc}
\usepackage{sectsty}
\usepackage{longtable}
\usepackage{array}
\usepackage{float}
\usepackage{fancyhdr}
\usepackage{enumitem}
\usepackage[colorlinks=true, allcolors=Azul-UCR]{hyperref}

%%%%%%%%%%%%%%%%%%%%%%%%%%%%%%%%%%%%%%%%%%%%%%%%%%
%%%%%%%%%% Colores UCR y estilos %%%%%%%%%%%%%%%%%%
%%%%%%%%%%%%%%%%%%%%%%%%%%%%%%%%%%%%%%%%%%%%%%%%%%
\definecolor{Celeste-UCR}{RGB}{0, 192, 243}
\definecolor{Azul-UCR}{RGB}{0, 93, 164}
\definecolor{Verde-UCR}{RGB}{109, 192, 103}
\definecolor{Amarillo-UCR}{RGB}{255, 224, 106}

\sectionfont{\color{Azul-UCR}}
\subsectionfont{\color{Azul-UCR}}
\subsubsectionfont{\color{Azul-UCR}}

\setlength{\parindent}{0pt}
\setlength{\parskip}{5pt}
\setlist{itemsep=2pt, topsep=3pt}
\setcounter{secnumdepth}{1}

\pagestyle{fancy}
\fancyhf{}
\renewcommand{\headrulewidth}{0.4pt}
\fancyhead[L]{\small\color{Azul-UCR} Resumen ejecutivo --- Bachillerato y Licenciatura en Matemática}
\fancyhead[R]{\small\color{Azul-UCR} Universidad de Costa Rica}
\fancyfoot[C]{\small\thepage}

\newcommand{\N}{\mathbb{N}}
\newcommand{\Z}{\mathbb{Z}}
\newcommand{\C}{\mathbb{C}}
\newcommand{\Q}{\mathbb{Q}}
\newcommand{\I}{\mathbb{I}}
\newcommand{\R}{\mathbb{R}}
\newcommand{\op}[1]{\operatorname{#1}}
\newcommand{\bb}[1]{\mathbb{#1}}
\newcommand{\fr}[1]{\mathfrak{#1}}
\newcommand{\legendre}[2]{\left( \dfrac{#1}{#2} \right)}
\newcommand{\mcd}{\operatorname{mcd}}
\newcommand{\MA}[1]{MA-#1}
\newcommand{\rojo}[1]{{#1}}

% Información que la propuesta no contiene y que debe completar la unidad académica.
\newcommand{\pendiente}[1]{%
  \par\smallskip\noindent\fcolorbox{Amarillo-UCR!90!black}{Amarillo-UCR!35}{%
    \parbox{\dimexpr\linewidth-2\fboxsep-2\fboxrule\relax}{\small\textbf{Por completar por la unidad académica:} #1}}\par\smallskip}

\newcommand{\encabezadotabla}[1]{\rowcolor{Azul-UCR}\textcolor{white}{\textbf{#1}}}

\title{Resumen ejecutivo de la propuesta curricular}
\author{Escuela de Matemática}
\date{\today}

\begin{document}

%%%%%%%%%%%%%%%%%%%%%%%%%%%%%%%%%%%%%%%%%%%%%%%%%%
%%%%%%%%%%%%%%%%%%%% Portada %%%%%%%%%%%%%%%%%%%%%
%%%%%%%%%%%%%%%%%%%%%%%%%%%%%%%%%%%%%%%%%%%%%%%%%%
\begin{titlepage}
\thispagestyle{empty}
\begin{tikzpicture}[remember picture, overlay]
  \fill[Azul-UCR] (current page.south west) rectangle (current page.north east);
  \fill[Celeste-UCR] (current page.south west) rectangle ([xshift=0.85cm]current page.north west);
  \fill[Amarillo-UCR] ([xshift=0.85cm]current page.south west) rectangle ([xshift=1.00cm]current page.north west);
  \fill[Celeste-UCR, opacity=0.16] ([xshift=-1.6cm, yshift=-1.4cm]current page.north east) circle (5.4cm);
  \fill[Verde-UCR, opacity=0.14] ([xshift=-0.4cm, yshift=-9.2cm]current page.north east) circle (3.6cm);
  \fill[Amarillo-UCR, opacity=0.10] ([xshift=1.8cm, yshift=3.4cm]current page.south east) circle (4.2cm);

  \begin{scope}[shift={([xshift=-7.55cm, yshift=-12.6cm]current page.north east)}]
    \foreach \x in {0,1,2,3} {
      \foreach \y in {0,1,2,3,4} {
        \draw[white, opacity=0.18, line width=0.6pt] (\x*1.42, -\y*1.28) -- (\x*1.42, -\y*1.28-1.28);
      }
    }
    \foreach \x in {0,1,2} {
      \foreach \y in {0,1,2,3,4,5} {
        \draw[white, opacity=0.18, line width=0.6pt] (\x*1.42, -\y*1.28) -- (\x*1.42+1.42, -\y*1.28);
      }
    }
    \foreach \x in {0,1,2,3} {
      \foreach \y in {0,1,2,3,4,5} {
        \fill[white, opacity=0.22] (\x*1.42, -\y*1.28) circle (0.19);
      }
    }
    \draw[Celeste-UCR, line width=1.55pt] (0,0) -- (1.42,-1.28) -- (2.84,-1.28) -- (2.84,-3.84) -- (4.26,-5.12);
    \fill[Amarillo-UCR] (0,0) circle (0.30);
    \fill[Celeste-UCR] (1.42,-1.28) circle (0.24);
    \fill[Celeste-UCR] (2.84,-1.28) circle (0.24);
    \fill[Celeste-UCR] (2.84,-3.84) circle (0.24);
    \fill[Verde-UCR] (4.26,-5.12) circle (0.30);
  \end{scope}

  \node[anchor=north west] at ([xshift=2.15cm, yshift=-1.45cm]current page.north west) {
    \includegraphics[height=1.45cm]{assets/logos/firma-ucr-blanco.png}
  };

  \node[anchor=north west, inner sep=0pt] at ([xshift=2.15cm, yshift=-4.15cm]current page.north west) {
    \begin{minipage}{12.4cm}
      \color{white}\raggedright
      {\color{Celeste-UCR}\fontsize{12}{14}\selectfont\bfseries RESUMEN EJECUTIVO \textbullet{} OPES--CONARE\par}
      \vspace{10pt}
      {\fontsize{36}{40}\selectfont\bfseries
      Resumen\\
      ejecutivo de la\\
      propuesta\\
      curricular\par}
      \vspace{12pt}
      {\color{Amarillo-UCR}\rule{3.15cm}{3.2pt}\par}
      \vspace{14pt}
      {\fontsize{16}{20}\selectfont
      Bachillerato y Licenciatura en Matemática \\ con énfasis en Matemática Pura y Aplicada\par}
    \end{minipage}
  };

  \node[anchor=south west, inner sep=0pt] at ([xshift=2.15cm, yshift=2.35cm]current page.south west) {
    \begin{minipage}{10cm}
      \color{white}\raggedright
      \includegraphics[height=1.7cm]{assets/logos/logo-mate.png}\\[12pt]
      {\fontsize{12}{15}\selectfont Escuela de Matemática, Universidad de Costa Rica\par}
      {\fontsize{12}{15}\selectfont Comisión de Estructura Curricular\par}
      {\color{Celeste-UCR}\fontsize{12}{15}\selectfont \the\year\par}
    \end{minipage}
  };
\end{tikzpicture}
\end{titlepage}

\addtocounter{page}{1}

Este documento resume la propuesta de modificación curricular de la carrera \textbf{Bachillerato y Licenciatura en Matemática con énfasis en Matemática Pura y Aplicada} de la Universidad de Costa Rica, y se organiza según los once apartados solicitados por la Oficina de Planificación de la Educación Superior (OPES) del CONARE. El detalle completo se encuentra en el documento \emph{Propuesta de modificación curricular} aprobado por la Asamblea de Escuela de Matemática el 28 de octubre de 2026.

%%%%%%%%%%%%%%%%%%%%%%%%%%%%%%%%%%%%%%%%%%%%%%%%%%
\section{Datos generales}

\begingroup
\renewcommand{\arraystretch}{1.35}
\arrayrulecolor{Azul-UCR!80}
\begin{longtable}{|>{\columncolor{Celeste-UCR!15}\raggedright\arraybackslash\bfseries}p{4.3cm}|>{\raggedright\arraybackslash}p{12.5cm}|}
\hline
Universidad & Universidad de Costa Rica (UCR). \\ \hline
Unidad académica responsable & Escuela de Matemática, Facultad de Ciencias. La carrera se imparte en la Sede Rodrigo Facio (Campus Universitario). \\ \hline
Otras unidades académicas participantes & Departamento de Ciencias Actuariales de la Escuela de Matemática (cursos con sigla CA) y las escuelas que ofrecen los cursos optativos de otra disciplina: Biología, Centroamericana de Geología, Ciencias de la Computación e Informática, Economía, Estadística, Filosofía, Física y Química. \\ \hline
Nombre de la carrera & Bachillerato y Licenciatura en Matemática con énfasis en Matemática Pura y Aplicada. \\ \hline
Grados y títulos &
\begin{minipage}[t]{12.2cm}
\begin{itemize}[leftmargin=1.2em, itemsep=1pt, topsep=2pt, parsep=0pt]
  \item Bachillerato en Matemática
  \item Licenciatura en Matemática, énfasis en Matemática Pura.
  \item Licenciatura en Matemática, énfasis en Matemática Aplicada.
\end{itemize}
\end{minipage} \\ \hline
Duración &
\begin{minipage}[t]{12.2cm}
\begin{itemize}[leftmargin=1.2em, itemsep=1pt, topsep=2pt, parsep=0pt]
  \item \textbf{Número de ciclos:} 8 ciclos para el bachillerato (4 años) y 2 ciclos adicionales para la licenciatura (1 año más); en total, 10 ciclos.
  \item \textbf{Semanas por ciclo:} 16 semanas lectivas.
  \item \textbf{Ciclos por año:} 2 (I y II ciclo lectivo).
\end{itemize}
\end{minipage} \\ \hline
Créditos &
Bachillerato: 130 créditos (énfasis en Matemática Pura) o 131 créditos (énfasis en Matemática Aplicada), de los cuales 50 corresponden al bloque común (ciclos I a III). Licenciatura: 30 créditos adicionales (ciclos IX y X). \\ \hline
Modalidad & Presencial. Algunas clases prácticas, de estilo laboratorio computacional, pueden ofrecerse de forma virtual parcial o total. \\ \hline
Plan que sustituye & Plan vigente 210401-2 (``plan 02''). El nuevo plan (``plan 2028'') rige para las personas que ingresen a partir de 2028. \\ \hline
\end{longtable}
\endgroup


%%%%%%%%%%%%%%%%%%%%%%%%%%%%%%%%%%%%%%%%%%%%%%%%%%
\section{Justificación}

La propuesta responde a la necesidad de contar con personas profesionales en matemática con excelencia académica, ética y sentido crítico, capaces de aplicar, investigar y comunicar su conocimiento de forma colaborativa. Atiende las demandas de investigación y docencia universitaria, de modelación matemática para otras disciplinas, de análisis de datos y de consultoría cuantitativa en el sector productivo, y contribuye al proyecto humanista de la Universidad mediante la generación de capital humano avanzado, la innovación en sectores estratégicos y la cultura científica del país.

La matemática ocupa un lugar central en el desarrollo científico, tecnológico y social contemporáneo. La digitalización de la información, el avance de la inteligencia artificial, el manejo de datos a gran escala y la modelación de fenómenos complejos han ampliado el objeto de estudio de la disciplina y los ámbitos en que el matemático puede incidir. La UNESCO proclamó en 2019 el 14 de marzo como Día Internacional de las Matemáticas y subraya su papel en el logro de los Objetivos de Desarrollo Sostenible y en la formación de capacidades para enfrentar desafíos como el cambio climático, la salud pública y la innovación tecnológica.

La demanda laboral respalda esta pertinencia. En Costa Rica, la empleabilidad de las carreras STEM fue de 95,8\,\% frente a 92,3\,\% en las no STEM, y carreras afines como Estadística o Ciencias Actuariales presentan niveles de desempleo muy bajos. A nivel internacional, la Oficina de Trabajo Estadístico de Estados Unidos proyectó un aumento del 28\,\% en los empleos de ocupaciones matemáticas entre 2016 y 2026, impulsado por el \emph{big data} y la analítica de datos.

Las personas profesionales en matemática encuentran oportunidades en instituciones de educación superior y de investigación, en organismos públicos y privados que requieren modelación y análisis cuantitativo, y en sectores como finanzas, banca, economía, tecnología y servicios basados en datos. 

Sobre el contexto universitario en que se desarrolla la carrera actual tenemos:

\begin{itemize}
  \item \textbf{Trayectoria de la carrera.} La Escuela de Matemática ofrece el Bachillerato y la Licenciatura desde 1974 (con antecedentes desde 1940), y desarrolla investigación y docencia de posgrado mediante la Maestría en Matemáticas, el CIMM, el CIMPA y la \emph{Revista de Matemática: Teoría y Aplicaciones}. Cuenta con 128 docentes, 41 de ellos con grado de doctor.
  \item \textbf{Actualización del plan.} El último rediseño integral del plan data de hace más de tres décadas. La propuesta se apoya en el perfil de egreso y los marcos referenciales elaborados entre 2018 y 2020 con apoyo del Centro de Evaluación Académica (CEA), mediante entrevistas, grupos focales, análisis estructurales de la actividad y consultas a la comunidad académica; el perfil fue aprobado en la Asamblea de Escuela en 2020.
  \item \textbf{Situación de la población estudiantil.} Entre 2014 y 2026 se graduaron 72 personas del bachillerato, con una duración real media de 5,3 años frente a los 4 años del plan; solo el 22,2\,\% concluyó en 4 años o menos. De las 543 solicitudes aceptadas entre 2012 y 2018, apenas el 9,9\,\% se había graduado en 2026. La propuesta busca atender el rezago mediante cursos preliminares en el primer ciclo y una introducción más gradual de la abstracción.
  \item \textbf{Admisión.} La oferta de cupos y la demanda han disminuido desde 2002; desde 2017 el cupo de primer ingreso es de 32 plazas, y la ocupación se recuperó a partir de 2024.
\end{itemize}

La creación de la carrera de Educación Matemática (2017) y la separación de las funciones actuariales hacia Ciencias Actuariales redefinen el espacio ocupacional de la persona matemática, orientándola hacia la investigación, la docencia universitaria, la modelación, el análisis de datos y el trabajo interdisciplinario.

La actualización de la formación de profesionales en matemática es fundamental para mantener el rol preponderante que la disciplina tiene en la formación de profesionales de otras carreras, la investigación en matemática y otras disciplinas así como en los programas de acción social donde son necesarios conocimientos de matemáticas.

La propuesta también busca favorecer el éxito académico de la población estudiantil, procurando un enfoque pedagógico más moderno y acorde con el perfil de la población de entrada. 



%%%%%%%%%%%%%%%%%%%%%%%%%%%%%%%%%%%%%%%%%%%%%%%%%%
\section{Objetivos de la carrera}

\textbf{Objetivo general.} Formar profesionales en matemática con excelencia académica, ética y sentido crítico, capaces de aplicar, investigar y comunicar su conocimiento de forma colaborativa, para contribuir al desarrollo productivo y social del país desde una perspectiva de bien común.

\textbf{Objetivos específicos.}
\begin{enumerate}
  \item Formar una persona profesional que impacte de manera positiva el entorno productivo, público y privado, mediante la aplicación del conocimiento matemático.
  \item Formar una persona profesional capaz de transmitir sus ideas de manera clara y respetuosa, que interactúe con profesionales de otras disciplinas y muestre gusto por compartir su conocimiento.
  \item Formar una persona profesional con un alto compromiso con la investigación, tanto teórica como aplicada.
  \item Formar una persona profesional con convicciones éticas y pensamiento crítico y reflexivo sobre su práctica profesional.
  \item Formar una persona profesional comprometida con el bien común y la justicia, que utilice su conocimiento para construir una sociedad más justa y solidaria.
  \item Incentivar una actitud positiva hacia el trabajo colaborativo con otros especialistas, en la búsqueda de soluciones que mejoren la calidad de vida y el entorno social.
  \item Formar una persona profesional dispuesta a aprender siempre, a buscar nuevas formas de hacer las cosas y a mantenerse actualizada en su área de especialidad.
\end{enumerate}

{\footnotesize Los objetivos corresponden a la Declaración de propósitos (propósito general y propósitos específicos) de la propuesta curricular.\par}

%%%%%%%%%%%%%%%%%%%%%%%%%%%%%%%%%%%%%%%%%%%%%%%%%%
\section{Perfil académico-profesional}

En la presente sección se presenta una síntesis del perfil de egreso.

%\newpage 

\subsection*{Bachillerato en matemática, énfasis en matemática pura}


\begin{table}[H]
\centering
\renewcommand{\arraystretch}{1.3}
\begin{tabular}{|c|p{9cm}|}
\hline
\cellcolor[HTML]{00C0F3}\textbf{Código(s)} & \cellcolor[HTML]{00C0F3}\textbf{Enunciado de la dimensión declarativa} \\
\hline
CDP01 & Conoce y comprende los fundamentos matemáticos —incluyendo la lógica, la teoría de números y la geometría euclídea— y los principios formales que sustentan el razonamiento y la argumentación rigurosa. \\
\hline
CDP02 & Conoce el uso de software especializado, los fundamentos de la programación y la computación científica, así como las técnicas básicas de análisis numérico para formular, explorar y resolver problemas matemáticos mediante enfoques computacionales. \\
\hline
CDP03 & Conoce los conceptos, métodos y aplicaciones básicos del cálculo y del álgebra lineal como fundamentos de la formación matemática. \\
\hline
CDP04 & Conoce los conceptos y técnicas de las ecuaciones diferenciales, la teoría de la medida, la probabilidad y los procesos estocásticos, así como del análisis complejo y funcional. \\
\hline
CDP05 & Conoce los fundamentos y resultados básicos de la geometría y la topología. \\
\hline
CDP06 & Conoce las estructuras y resultados fundamentales del álgebra abstracta. 
\\
\hline
CDP07 & Conoce las estructuras y resultados fundamentales de algún área de especialización en matemática pura más allá de las antes expuestas. 
\\
\hline
\end{tabular}
\caption{Conocimientos declarativos o saber conocer.}
\end{table}

%\newpage

\begin{table}[H]
\centering
\renewcommand{\arraystretch}{1.3}
\begin{tabular}{|c|p{9cm}|}
\hline
\cellcolor[HTML]{00C0F3}\textbf{Código(s)}  & \cellcolor[HTML]{00C0F3}\textbf{Enunciado de la dimensión procedimental} \\
\hline
CPP01 & Aplica pensamiento operacional, abstracto y formal para analizar, resolver y sintetizar problemas matemáticos. \\
\hline
CPP02 & Formula conjeturas, resuelve problemas, interpreta resultados y se comunica de manera coherente en lenguaje matemático formal. \\
\hline
CPP03 & Utiliza tecnologías y software científico para explorar, validar y resolver problemas matemáticos.\\
\hline
CPP04 & Participa activamente en discusiones científicas y divulga el conocimiento matemático de forma accesible a diferentes audiencias.\\
\hline
CPP05 & Formula y demuestra rigurosamente resultados matemáticos, generalizando estructuras y construyendo argumentos lógicos válidos dentro de un sistema axiomático.\\
\hline
\end{tabular}
\caption{Conocimientos procedimentales o saber hacer.}
\end{table}

\begin{table}[H]
\centering
\renewcommand{\arraystretch}{1.3}
\begin{tabular}{|c|p{9cm}|}
\hline
\cellcolor[HTML]{00C0F3}\textbf{Código(s)} & \cellcolor[HTML]{00C0F3}\textbf{Enunciado de la dimensión actitudinal} \\
\hline
CAP01 & Manifiesta respeto por la diversidad humana, matemática, científica y cultural, fomentando un ambiente inclusivo y de valoración mutua. \\
\hline
CAP02 & Se comunica con claridad, colabora en equipos académicos y profesionales, y mantiene una disposición al autoaprendizaje y al liderazgo.\\
\hline
CAP03 & Demuestra conciencia social y participación ciudadana, con una visión sociocultural e integral de la sociedad en la que la ciencia y la matemática ocupan un lugar relevante.\\
\hline
CAP04 & Ejerce pensamiento crítico, actúa con integridad ética y evalúa con rigor las prácticas científicas y académicas.\\
\hline
CAP05 & Utiliza su conocimiento matemático para resolver problemas académicos y profesionales, adaptándolo a distintos contextos y necesidades.\\
\hline
CAP06 & Mantiene un interés permanente por los desarrollos de la investigación científica y matemática, valorando la actualización como parte de su formación profesional.\\
\hline
CAP07 & Muestra paciencia y perseverancia ante la abstracción sostenida, valorando el rigor y la elegancia argumentativa como fines en sí mismos, incluso cuando la solución de un problema no es inmediata.\\
\hline
\end{tabular}
\caption{Conocimientos actitudinales o saber ser.}
\end{table}


\subsection*{Bachillerato en matemática, énfasis en matemática aplicada}



\begin{table}[H]
\centering
\renewcommand{\arraystretch}{1.3}
\begin{tabular}{|c|p{9cm}|}
\hline
\cellcolor[HTML]{00C0F3}\textbf{Código(s)} & \cellcolor[HTML]{00C0F3}\textbf{Enunciado de la dimensión declarativa} \\
\hline
CDA01 & Conoce y comprende los fundamentos matemáticos —incluyendo la lógica, la teoría de números y la geometría euclídea— y los principios formales que sustentan el razonamiento y la argumentación rigurosa. \\
\hline
CDA02 & Conoce el uso de software especializado, los fundamentos de la programación y la computación científica, así como las técnicas básicas de análisis numérico para formular, explorar y resolver problemas matemáticos mediante enfoques computacionales. \\
\hline
CDA03 & Conoce los conceptos, métodos y aplicaciones básicos del cálculo y del álgebra lineal como fundamentos de la formación matemática. \\
\hline
CDA04 & Conoce los conceptos y técnicas de las ecuaciones diferenciales, la teoría de la medida, la probabilidad y los procesos estocásticos, así como del análisis complejo y funcional. \\
\hline
CDA05 & Conoce los principios de la matemática aplicada, por ejemplo, por medio de la modelación, la estadística o el análisis de datos, para resolver problemas y tomar decisiones en ámbitos como la industria, la banca y otros sectores.
\\
\hline
CDA06 & Conoce las principales tecnologías de aplicación de la matemática, por ejemplo, la programación y computación científica para aplicar conceptos teóricos a problemas reales.
\\
\hline
CDA07 & Conoce las estructuras y resultados fundamentales de algún área de especialización en matemática aplicada más allá de las antes expuestas. 
\\
\hline
\end{tabular}
\caption{Conocimientos declarativos o saber conocer.}
\end{table}

%\newpage

\begin{table}[H]
\centering
\renewcommand{\arraystretch}{1.3}
\begin{tabular}{|c|p{9cm}|}
\hline
\cellcolor[HTML]{00C0F3}\textbf{Código(s)}  & \cellcolor[HTML]{00C0F3}\textbf{Enunciado de la dimensión procedimental} \\
\hline
CPA01 & Aplica pensamiento operacional, abstracto y formal para analizar, resolver y sintetizar problemas matemáticos. \\
\hline
CPA02 & Formula conjeturas, resuelve problemas, interpreta resultados y se comunica de manera coherente en lenguaje matemático formal. \\
\hline
CPA03 & Utiliza tecnologías y software científico para explorar, validar y resolver problemas matemáticos.\\
\hline
CPA04 & Participa activamente en discusiones científicas y divulga el conocimiento matemático de forma accesible a diferentes audiencias.\\
\hline
CPA05 & Formula modelos matemáticos a partir de problemas del entorno, los resuelve mediante herramientas analíticas o computacionales, e interpreta y valida los resultados en el contexto original.\\
\hline
\end{tabular}
\caption{Conocimientos procedimentales o saber hacer.}
\end{table}


%\newpage

\begin{table}[H]
\centering
\renewcommand{\arraystretch}{1.3}
\begin{tabular}{|c|p{9cm}|}
\hline
\cellcolor[HTML]{00C0F3}\textbf{Código(s)} & \cellcolor[HTML]{00C0F3}\textbf{Enunciado de la dimensión actitudinal} \\
\hline
CAA01 & Manifiesta respeto por la diversidad humana, matemática, científica y cultural, fomentando un ambiente inclusivo y de valoración mutua. \\
\hline
CAA02 & Se comunica con claridad, colabora en equipos académicos y profesionales, y mantiene una disposición al autoaprendizaje y al liderazgo.\\
\hline
CAA03 & Demuestra conciencia social y participación ciudadana, con una visión sociocultural e integral de la sociedad en la que la ciencia y la matemática ocupan un lugar relevante.\\
\hline
CAA04 & Ejerce pensamiento crítico, actúa con integridad ética y evalúa con rigor las prácticas científicas y académicas.\\
\hline
CAA05 & Utiliza su conocimiento matemático para resolver problemas académicos y profesionales, adaptándolo a distintos contextos y necesidades.\\
\hline
CAA06 & Mantiene un interés permanente por los desarrollos de la investigación científica y matemática, valorando la actualización como parte de su formación profesional.\\
\hline
CAA07 & Asume con responsabilidad ética las implicaciones prácticas de los modelos y datos que utiliza, y muestra apertura al diálogo interdisciplinario y tolerancia a la incertidumbre propia de los fenómenos reales.\\
\hline
\end{tabular}
\caption{Conocimientos actitudinales o saber ser.}
\end{table}


\subsection*{Licenciatura en matemática, énfasis en matemática pura}


En adición a los conocimientos enunciados para el bachillerato, con la licenciatura se tendrían los siguientes.

\begin{table}[H]
\centering
\renewcommand{\arraystretch}{1.3}
\begin{tabular}{|c|p{9cm}|}
\hline
\cellcolor[HTML]{00C0F3}\textbf{Código(s)} & \cellcolor[HTML]{00C0F3}\textbf{Enunciado de la dimensión declarativa} \\
\hline
CDP08 & Conoce estructuras y resultados de nivel intermedio para algunas áreas medulares de especialización en matemática pura. 
\\
\hline
\end{tabular}
\caption{Conocimientos declarativos o saber conocer.}
\end{table}

%\newpage

\begin{table}[H]
\centering
\renewcommand{\arraystretch}{1.3}
\begin{tabular}{|c|p{9cm}|}
\hline
\cellcolor[HTML]{00C0F3}\textbf{Código(s)}  & \cellcolor[HTML]{00C0F3}\textbf{Enunciado de la dimensión procedimental} \\
\hline
CP05 & Aplica pensamiento operacional, abstracto y formal, con fuentes transversales de varias áreas, para formular, analizar y resolver problemas matemáticos. \\
\hline
\end{tabular}
\caption{Conocimientos procedimentales o saber hacer.}
\end{table}


\begin{table}[H]
\centering
\renewcommand{\arraystretch}{1.3}
\begin{tabular}{|c|p{9cm}|}
\hline
\cellcolor[HTML]{00C0F3}\textbf{Código(s)} & \cellcolor[HTML]{00C0F3}\textbf{Enunciado de la dimensión actitudinal} \\
\hline
CA06 & Muestra iniciativa y autonomía para formular preguntas, explorar problemas matemáticos y profundizar en ellos más allá de los procedimientos conocidos.\\
\hline
\end{tabular}
\caption{Conocimientos actitudinales o saber ser.}
\end{table}

\subsection*{Licenciatura en matemática, énfasis en matemática aplicada}

En adición a los conocimientos enunciados para el bachillerato, con la licenciatura se tendrían los siguientes.

\begin{table}[H]
\centering
\renewcommand{\arraystretch}{1.3}
\begin{tabular}{|c|p{9cm}|}
\hline
\cellcolor[HTML]{00C0F3}\textbf{Código(s)} & \cellcolor[HTML]{00C0F3}\textbf{Enunciado de la dimensión declarativa} \\
\hline
CDA08 & Conoce estructuras y resultados de nivel intermedio para algunas áreas medulares de especialización en matemática aplicada. 
\\
\hline
\end{tabular}
\caption{Conocimientos declarativos o saber conocer.}
\end{table}

%\newpage

\begin{table}[H]
\centering
\renewcommand{\arraystretch}{1.3}
\begin{tabular}{|c|p{9cm}|}
\hline
\cellcolor[HTML]{00C0F3}\textbf{Código(s)}  & \cellcolor[HTML]{00C0F3}\textbf{Enunciado de la dimensión procedimental} \\
\hline
CP05 & Aplica pensamiento operacional, abstracto y formal, con fuentes transversales de varias áreas, para formular, analizar y resolver problemas matemáticos. \\
\hline
\end{tabular}
\caption{Conocimientos procedimentales o saber hacer.}
\end{table}


\begin{table}[H]
\centering
\renewcommand{\arraystretch}{1.3}
\begin{tabular}{|c|p{9cm}|}
\hline
\cellcolor[HTML]{00C0F3}\textbf{Código(s)} & \cellcolor[HTML]{00C0F3}\textbf{Enunciado de la dimensión actitudinal} \\
\hline
CA06 & Muestra iniciativa y autonomía para formular preguntas, explorar problemas matemáticos y profundizar en ellos más allá de los procedimientos conocidos.\\
\hline
\end{tabular}
\caption{Conocimientos actitudinales o saber ser.}
\end{table}

%%%%%%%%%%%%%%%%%%%%%%%%%%%%%%%%%%%%%%%%%%%%%%%%%%
\section{Campo de inserción profesional}

Los principales campos de inserción identificados para las personas graduadas son:
\begin{itemize}
  \item \textbf{Investigación académica} (teórica y aplicada), en universidades, centros de investigación y redes científicas nacionales e internacionales.
  \item \textbf{Docencia universitaria}, en programas de pregrado y posgrado en matemática y disciplinas afines.
  \item \textbf{Modelación matemática}, en ingeniería, física, biología, epidemiología, ciencias ambientales y otras áreas.
  \item \textbf{Análisis de datos y ciencia de datos}, en empresas, instituciones financieras, organismos públicos y proyectos de investigación interdisciplinaria.
  \item \textbf{Sector productivo e industrial}, en consultoría cuantitativa, optimización, criptografía, simulación y apoyo a la innovación tecnológica.
  \item \textbf{Trabajo interdisciplinario}, como puente entre equipos de distintas disciplinas en proyectos científicos, tecnológicos o de política pública.
\end{itemize}

%%%%%%%%%%%%%%%%%%%%%%%%%%%%%%%%%%%%%%%%%%%%%%%%%%
\section{Requisitos de ingreso}

El ingreso a la carrera se realiza por las modalidades de admisión de la Universidad de Costa Rica y se rige por los procedimientos institucionales de admisión para carreras de grado.
\begin{itemize}
  \item \textbf{Primer ingreso}, mediante la nota de admisión que incorpora la Prueba de Aptitud Académica, con 32 cupos ordinarios y 3 cupos de primer ingreso diferido por año.
  \item \textbf{Traslado}, en tres modalidades: por nota de admisión, por rendimiento académico y por excelencia académica.
\end{itemize}

A partir de 2028, toda persona de primer ingreso, traslado o segunda carrera se incorpora directamente al nuevo plan; no se ofrecen nuevas admisiones al plan vigente. Las cohortes que ya cursan el plan 02 pueden pasar al nuevo plan mediante el derecho de opción y la tabla de convalidaciones.

No se identifican cursos nivelatorios específicos definidos para la presente propuesta. La permanencia y avance en el plan se regulan por los requisitos, correquisitos y normativa académica institucional aplicable.

%%%%%%%%%%%%%%%%%%%%%%%%%%%%%%%%%%%%%%%%%%%%%%%%%%
\section{Requisitos de graduación}

\subsection*{Bachillerato en Matemática (ambos énfasis)}
\begin{itemize}
  \item Aprobar el plan de estudios de los ciclos I a VIII: 130 créditos en el énfasis en Matemática Pura o 131 créditos en el énfasis en Matemática Aplicada. Esto incluye la Práctica Profesional del último ciclo (modalidad académica o profesional).
  \item Asistir, a lo largo de la carrera, a diez (10) actividades académicas válidas (charlas del Coloquio del Departamento de Matemática, seminarios, conferencias, ciclos especiales u otras actividades que disponga la Dirección del Departamento).
  \item Cumplir los requisitos que establece la normativa universitaria, tales como el Trabajo Comunal Universitario (TCU) y el Trabajo Final de Graduación (TFG).
\end{itemize}

\subsection*{Licenciatura en Matemática (ambos énfasis)}
\begin{itemize}
  \item Aprobar los 30 créditos del plan de licenciatura (ciclos IX y X): dos seminarios de estudios dirigidos del énfasis y cursos optativos.
  \item Asistir a cinco (5) actividades académicas válidas, distintas de las empleadas para el bachillerato.
  \item Cumplir los requisitos de la normativa universitaria para la obtención del grado de licenciatura.
\end{itemize}

%%%%%%%%%%%%%%%%%%%%%%%%%%%%%%%%%%%%%%%%%%%%%%%%%%
\section{Listado de cursos}

Las tablas siguientes presentan el nombre, los créditos y el ciclo de cada curso de la estructura curricular. Los cursos con sigla MA son propios de la Escuela de Matemática y los de sigla CA corresponden al Departamento de Ciencias Actuariales; los cursos optativos de otra disciplina (OPT-) los ofrecen otras unidades académicas de la Universidad.

\setcounter{table}{0}
% Generado por scripts/generate_resumen_ejecutivo.py. No editar a mano.

\subsection*{Cursos de la estructura curricular}

\begingroup
\small
\arrayrulecolor{Azul-UCR!80}%
\renewcommand{\arraystretch}{1.2}%
\rowcolors{2}{white}{Celeste-UCR!12}%
\begin{longtable}{|c|l|>{\raggedright\arraybackslash}p{8.6cm}|c|}
\caption{Bloque común (ciclos I a III)}\label{res:malla-comun}\\
\hline
\rowcolor{Azul-UCR}
\textcolor{white}{\textbf{Ciclo}} & \textcolor{white}{\textbf{Sigla}} & \textcolor{white}{\textbf{Nombre del curso}} & \textcolor{white}{\textbf{Créditos}} \\ \hline
\endfirsthead
\hline
\rowcolor{Azul-UCR}
\textcolor{white}{\textbf{Ciclo}} & \textcolor{white}{\textbf{Sigla}} & \textcolor{white}{\textbf{Nombre del curso}} & \textcolor{white}{\textbf{Créditos}} \\ \hline
\endhead
\hline
\endfoot
1 & EF- & Actividad Deportiva & 0 \\ \hline
1 & EG-I & Curso Integrado de Humanidades I & 6 \\ \hline
1 & MA-0151 & Fundamentos de Álgebra, Trigonometría y Geometría Analítica & 3 \\ \hline
1 & MA-0152 & Matemática Exploratoria & 3 \\ \hline
1 & RP- & Repertorio & 3 \\ \hline
\rowcolor{Verde-UCR!28}\multicolumn{4}{|l|}{\textit{Créditos del ciclo 1: 15}} \\ \hline
2 & EG- & Curso de Arte & 2 \\ \hline
2 & EG-II & Curso Integrado de Humanidades II & 6 \\ \hline
2 & LM-3039 & Inglés para Estadística I & 3 \\ \hline
2 & MA-0251 & Cálculo en una Variable & 4 \\ \hline
2 & MA-0252 & Introducción a las Demostraciones & 3 \\ \hline
\rowcolor{Verde-UCR!28}\multicolumn{4}{|l|}{\textit{Créditos del ciclo 2: 18}} \\ \hline
3 & SR-I & Seminario de Realidad Nacional I & 2 \\ \hline
3 & MA-0361 & Álgebra Lineal I & 4 \\ \hline
3 & LM-3040 & Inglés para Estadística II & 3 \\ \hline
3 & MA-0351 & Cálculo en Varias Variables & 4 \\ \hline
3 & MA-0341 & Matemática Computacional & 4 \\ \hline
\rowcolor{Verde-UCR!28}\multicolumn{4}{|l|}{\textit{Créditos del ciclo 3: 17}} \\ \hline
\rowcolor{Amarillo-UCR!45}\multicolumn{4}{|r|}{\textbf{Total créditos bloque común: 50}} \\ \hline
\end{longtable}
\endgroup

\begingroup
\small
\arrayrulecolor{Azul-UCR!80}%
\renewcommand{\arraystretch}{1.2}%
\rowcolors{2}{white}{Celeste-UCR!12}%
\begin{longtable}{|c|l|>{\raggedright\arraybackslash}p{8.6cm}|c|}
\caption{Énfasis en Matemática pura (ciclos IV a VIII)}\label{res:malla-pura}\\
\hline
\rowcolor{Azul-UCR}
\textcolor{white}{\textbf{Ciclo}} & \textcolor{white}{\textbf{Sigla}} & \textcolor{white}{\textbf{Nombre del curso}} & \textcolor{white}{\textbf{Créditos}} \\ \hline
\endfirsthead
\hline
\rowcolor{Azul-UCR}
\textcolor{white}{\textbf{Ciclo}} & \textcolor{white}{\textbf{Sigla}} & \textcolor{white}{\textbf{Nombre del curso}} & \textcolor{white}{\textbf{Créditos}} \\ \hline
\endhead
\hline
\endfoot
4 & SR-II & Seminario de Realidad Nacional II & 2 \\ \hline
4 & MA-0461 & Álgebra Lineal II & 4 \\ \hline
4 & MA-0451 & Principios de Análisis I & 4 \\ \hline
4 & MA-0471 & Introducción a la Geometría Diferencial & 3 \\ \hline
4 & MA-0496 & Teoría de Números & 3 \\ \hline
\rowcolor{Verde-UCR!28}\multicolumn{4}{|l|}{\textit{Créditos del ciclo 4: 16}} \\ \hline
5 & MA-0571 & Álgebra Abstracta I & 4 \\ \hline
5 & MA-0555 & Principios de Análisis II & 4 \\ \hline
5 & MA-0556 & Ecuaciones Diferenciales & 4 \\ \hline
5 & OPT- & Optativa de Otra Disciplina & 3 \\ \hline
\rowcolor{Verde-UCR!28}\multicolumn{4}{|l|}{\textit{Créditos del ciclo 5: 15}} \\ \hline
6 & MA-0661 & Álgebra Abstracta II & 5 \\ \hline
6 & MA-0635 & Topología General & 4 \\ \hline
6 & MA-0501 & Análisis Numérico I & 4 \\ \hline
6 & OPT- & Optativa de Otra Disciplina & 3 \\ \hline
\rowcolor{Verde-UCR!28}\multicolumn{4}{|l|}{\textit{Créditos del ciclo 6: 16}} \\ \hline
7 & MA-0730 & Análisis Real I & 5 \\ \hline
7 & MA-0735 & Análisis Complejo & 4 \\ \hline
7 & MA-0880 & Comunicación en las Ciencias & 2 \\ \hline
7 & OPT-001 & Optativa Núcleo en Matemática Pura & 4 \\ \hline
\rowcolor{Verde-UCR!28}\multicolumn{4}{|l|}{\textit{Créditos del ciclo 7: 15}} \\ \hline
8 & MA-0881 & Práctica Profesional en Matemática Pura & 6 \\ \hline
8 & OPT-001 & Optativa Núcleo en Matemática Pura & 8 \\ \hline
8 & OPT-003 & Optativa General en Matemática & 4 \\ \hline
\rowcolor{Verde-UCR!28}\multicolumn{4}{|l|}{\textit{Créditos del ciclo 8: 18}} \\ \hline
\rowcolor{Amarillo-UCR!45}\multicolumn{4}{|r|}{\textbf{Total créditos bachillerato pura: 130}} \\ \hline
\end{longtable}
\endgroup

\begingroup
\small
\arrayrulecolor{Azul-UCR!80}%
\renewcommand{\arraystretch}{1.2}%
\rowcolors{2}{white}{Celeste-UCR!12}%
\begin{longtable}{|c|l|>{\raggedright\arraybackslash}p{8.6cm}|c|}
\caption{Énfasis en Matemática aplicada (ciclos IV a VIII)}\label{res:malla-aplicada}\\
\hline
\rowcolor{Azul-UCR}
\textcolor{white}{\textbf{Ciclo}} & \textcolor{white}{\textbf{Sigla}} & \textcolor{white}{\textbf{Nombre del curso}} & \textcolor{white}{\textbf{Créditos}} \\ \hline
\endfirsthead
\hline
\rowcolor{Azul-UCR}
\textcolor{white}{\textbf{Ciclo}} & \textcolor{white}{\textbf{Sigla}} & \textcolor{white}{\textbf{Nombre del curso}} & \textcolor{white}{\textbf{Créditos}} \\ \hline
\endhead
\hline
\endfoot
4 & MA-0461 & Álgebra Lineal II & 4 \\ \hline
4 & MA-0451 & Principios de Análisis I & 4 \\ \hline
4 & CA-0721 & Probabilidad & 4 \\ \hline
4 & CA-0204 & Herramientas de Ciencia de Datos I & 4 \\ \hline
\rowcolor{Verde-UCR!28}\multicolumn{4}{|l|}{\textit{Créditos del ciclo 4: 16}} \\ \hline
5 & SR-II & Seminario de Realidad Nacional II & 2 \\ \hline
5 & MA-0555 & Principios de Análisis II & 4 \\ \hline
5 & MA-0556 & Ecuaciones Diferenciales & 4 \\ \hline
5 & CA-0303 & Estadística Actuarial I & 4 \\ \hline
5 & OPT- & Optativa de Otra Disciplina & 3 \\ \hline
\rowcolor{Verde-UCR!28}\multicolumn{4}{|l|}{\textit{Créditos del ciclo 5: 17}} \\ \hline
6 & MA-0501 & Análisis Numérico I & 4 \\ \hline
6 & MA-0635 & Topología General & 4 \\ \hline
6 & MA-0647 & Modelación Matemática & 4 \\ \hline
6 & OPT- & Optativa de Otra Disciplina & 3 \\ \hline
\rowcolor{Verde-UCR!28}\multicolumn{4}{|l|}{\textit{Créditos del ciclo 6: 15}} \\ \hline
7 & MA-0730 & Análisis Real I & 5 \\ \hline
7 & MA-0735 & Análisis Complejo & 4 \\ \hline
7 & CA-0411 & Análisis de Datos I & 4 \\ \hline
7 & MA-0880 & Comunicación en las Ciencias & 2 \\ \hline
\rowcolor{Verde-UCR!28}\multicolumn{4}{|l|}{\textit{Créditos del ciclo 7: 15}} \\ \hline
8 & MA-0883 & Práctica Profesional en Matemática Aplicada & 6 \\ \hline
8 & OPT-002 & Optativa Núcleo en Matemática Aplicada & 8 \\ \hline
8 & OPT-003 & Optativa General en Matemática & 4 \\ \hline
\rowcolor{Verde-UCR!28}\multicolumn{4}{|l|}{\textit{Créditos del ciclo 8: 18}} \\ \hline
\rowcolor{Amarillo-UCR!45}\multicolumn{4}{|r|}{\textbf{Total créditos bachillerato aplicada: 131}} \\ \hline
\end{longtable}
\endgroup

\begingroup
\small
\arrayrulecolor{Azul-UCR!80}%
\renewcommand{\arraystretch}{1.2}%
\rowcolors{2}{white}{Celeste-UCR!12}%
\begin{longtable}{|c|l|>{\raggedright\arraybackslash}p{8.6cm}|c|}
\caption{Énfasis en Matemática pura (licenciatura, ciclos IX y X)}\label{res:lic-pura}\\
\hline
\rowcolor{Azul-UCR}
\textcolor{white}{\textbf{Ciclo}} & \textcolor{white}{\textbf{Sigla}} & \textcolor{white}{\textbf{Nombre del curso}} & \textcolor{white}{\textbf{Créditos}} \\ \hline
\endfirsthead
\hline
\rowcolor{Azul-UCR}
\textcolor{white}{\textbf{Ciclo}} & \textcolor{white}{\textbf{Sigla}} & \textcolor{white}{\textbf{Nombre del curso}} & \textcolor{white}{\textbf{Créditos}} \\ \hline
\endhead
\hline
\endfoot
9 & MA-0982 & Seminario de Estudios Dirigidos en Matemática Pura I & 7 \\ \hline
9 & OPT-001 & Optativa Núcleo en Matemática Pura & 4 \\ \hline
9 & OPT-003 & Optativa General en Matemática & 4 \\ \hline
\rowcolor{Verde-UCR!28}\multicolumn{4}{|l|}{\textit{Créditos del ciclo 9: 15}} \\ \hline
10 & MA-1081 & Seminario de Estudios Dirigidos en Matemática Pura II & 7 \\ \hline
10 & OPT-003 & Optativa General en Matemática & 8 \\ \hline
\rowcolor{Verde-UCR!28}\multicolumn{4}{|l|}{\textit{Créditos del ciclo 10: 15}} \\ \hline
\rowcolor{Amarillo-UCR!45}\multicolumn{4}{|r|}{\textbf{Total créditos licenciatura pura: 30}} \\ \hline
\end{longtable}
\endgroup

\begingroup
\small
\arrayrulecolor{Azul-UCR!80}%
\renewcommand{\arraystretch}{1.2}%
\rowcolors{2}{white}{Celeste-UCR!12}%
\begin{longtable}{|c|l|>{\raggedright\arraybackslash}p{8.6cm}|c|}
\caption{Énfasis en Matemática aplicada (licenciatura, ciclos IX y X)}\label{res:lic-aplicada}\\
\hline
\rowcolor{Azul-UCR}
\textcolor{white}{\textbf{Ciclo}} & \textcolor{white}{\textbf{Sigla}} & \textcolor{white}{\textbf{Nombre del curso}} & \textcolor{white}{\textbf{Créditos}} \\ \hline
\endfirsthead
\hline
\rowcolor{Azul-UCR}
\textcolor{white}{\textbf{Ciclo}} & \textcolor{white}{\textbf{Sigla}} & \textcolor{white}{\textbf{Nombre del curso}} & \textcolor{white}{\textbf{Créditos}} \\ \hline
\endhead
\hline
\endfoot
9 & MA-0984 & Seminario de Estudios Dirigidos en Matemática Aplicada I & 7 \\ \hline
9 & OPT-002 & Optativa Núcleo en Matemática Aplicada & 4 \\ \hline
9 & OPT-003 & Optativa General en Matemática & 4 \\ \hline
\rowcolor{Verde-UCR!28}\multicolumn{4}{|l|}{\textit{Créditos del ciclo 9: 15}} \\ \hline
10 & MA-1082 & Seminario de Estudios Dirigidos en Matemática Aplicada II & 7 \\ \hline
10 & OPT-003 & Optativa General en Matemática & 8 \\ \hline
\rowcolor{Verde-UCR!28}\multicolumn{4}{|l|}{\textit{Créditos del ciclo 10: 15}} \\ \hline
\rowcolor{Amarillo-UCR!45}\multicolumn{4}{|r|}{\textbf{Total créditos licenciatura aplicada: 30}} \\ \hline
\end{longtable}
\endgroup

\subsection*{Cursos optativos de Matemática}

\begingroup
\small
\arrayrulecolor{Azul-UCR!80}%
\renewcommand{\arraystretch}{1.2}%
\rowcolors{2}{white}{Celeste-UCR!12}%
\begin{longtable}{|l|>{\raggedright\arraybackslash}p{9.4cm}|c|}
\caption{Optativas de núcleo en Matemática Pura (OPT-001)}\label{res:opt001}\\
\hline
\rowcolor{Azul-UCR}
\textcolor{white}{\textbf{Sigla}} & \textcolor{white}{\textbf{Nombre del curso}} & \textcolor{white}{\textbf{Créditos}} \\ \hline
\endfirsthead
\hline
\rowcolor{Azul-UCR}
\textcolor{white}{\textbf{Sigla}} & \textcolor{white}{\textbf{Nombre del curso}} & \textcolor{white}{\textbf{Créditos}} \\ \hline
\endhead
\hline
\endfoot
MA-0526 & Combinatoria & 4 \\ \hline
MA-0528 & Sistemas Dinámicos & 4 \\ \hline
MA-0838 & Geometría Diferencial & 4 \\ \hline
MA-0837 & Geometría Algebraica & 4 \\ \hline
MA-0734 & Optimización & 4 \\ \hline
MA-0756 & Ecuaciones Diferenciales Parciales & 5 \\ \hline
MA-0781 & Lógica & 5 \\ \hline
MA-0783 & Integración & 4 \\ \hline
MA-0788 & Análisis Real II & 4 \\ \hline
MA-0835 & Topología Algebraica & 4 \\ \hline
MA-0834 & Análisis Funcional & 4 \\ \hline
MA-0863 & Teoría Analítica de Números & 4 \\ \hline
MA-0862 & Análisis Armónico & 4 \\ \hline
MA-0861 & Teoría de Modelos & 4 \\ \hline
MA-0841 & Probabilidad & 4 \\ \hline
MA-0849 & Teoría de Módulos & 5 \\ \hline
MA-0858 & Teoría de Conjuntos & 4 \\ \hline
MA-0876 & Teoría Algebraica de Números & 4 \\ \hline
MA-0889 & Álgebra Conmutativa & 4 \\ \hline
MA-0920 & Ecuaciones en Derivadas Parciales Numéricas & 4 \\ \hline
\end{longtable}
\endgroup

\begingroup
\small
\arrayrulecolor{Azul-UCR!80}%
\renewcommand{\arraystretch}{1.2}%
\rowcolors{2}{white}{Celeste-UCR!12}%
\begin{longtable}{|l|>{\raggedright\arraybackslash}p{9.4cm}|c|}
\caption{Optativas de núcleo en Matemática Aplicada (OPT-002)}\label{res:opt002}\\
\hline
\rowcolor{Azul-UCR}
\textcolor{white}{\textbf{Sigla}} & \textcolor{white}{\textbf{Nombre del curso}} & \textcolor{white}{\textbf{Créditos}} \\ \hline
\endfirsthead
\hline
\rowcolor{Azul-UCR}
\textcolor{white}{\textbf{Sigla}} & \textcolor{white}{\textbf{Nombre del curso}} & \textcolor{white}{\textbf{Créditos}} \\ \hline
\endhead
\hline
\endfoot
CA-0304 & Herramientas de Ciencia de Datos II & --- \\ \hline
CA-0307 & Estadística Actuarial II & --- \\ \hline
CA-0414 & Modelos Lineales y de Sobrevivencia & --- \\ \hline
CA-0512 & Análisis de Datos II & --- \\ \hline
CA-0602 & Introducción a la Optimización & --- \\ \hline
CA-0612 & Series de Tiempo & --- \\ \hline
MA-0477 & Ecuaciones Diferenciales Parciales Aplicadas & 4 \\ \hline
\end{longtable}
\endgroup

\begingroup
\small
\arrayrulecolor{Azul-UCR!80}%
\renewcommand{\arraystretch}{1.2}%
\rowcolors{2}{white}{Celeste-UCR!12}%
\begin{longtable}{|l|>{\raggedright\arraybackslash}p{9.4cm}|c|}
\caption{Optativas generales en Matemática (OPT-003)}\label{res:opt003}\\
\hline
\rowcolor{Azul-UCR}
\textcolor{white}{\textbf{Sigla}} & \textcolor{white}{\textbf{Nombre del curso}} & \textcolor{white}{\textbf{Créditos}} \\ \hline
\endfirsthead
\hline
\rowcolor{Azul-UCR}
\textcolor{white}{\textbf{Sigla}} & \textcolor{white}{\textbf{Nombre del curso}} & \textcolor{white}{\textbf{Créditos}} \\ \hline
\endhead
\hline
\endfoot
CA-0304 & Herramientas de Ciencia de Datos II & --- \\ \hline
CA-0303 & Estadística Actuarial I & 4 \\ \hline
CA-0403 & Estadística Actuarial II & --- \\ \hline
CA-0414 & Modelos Lineales y de Sobrevivencia & --- \\ \hline
CA-0512 & Análisis de Datos II & --- \\ \hline
CA-0602 & Introducción a la Optimización & --- \\ \hline
CA-0612 & Series de Tiempo & --- \\ \hline
MA-0477 & Ecuaciones Diferenciales Parciales Aplicadas & 4 \\ \hline
MA-0526 & Combinatoria & 4 \\ \hline
MA-0528 & Sistemas Dinámicos & 4 \\ \hline
MA-0838 & Geometría Diferencial & 4 \\ \hline
MA-0837 & Geometría Algebraica & 4 \\ \hline
MA-0734 & Optimización & 4 \\ \hline
MA-0756 & Ecuaciones Diferenciales Parciales & 5 \\ \hline
MA-0759 & Tópicos de Matemática Financiera & 4 \\ \hline
MA-0781 & Lógica & 5 \\ \hline
MA-0783 & Integración & 4 \\ \hline
MA-0788 & Análisis Real II & 4 \\ \hline
MA-0790 & Tópicos de Análisis & 4 \\ \hline
MA-0791 & Tópicos de Probabilidad & 4 \\ \hline
MA-0792 & Tópicos de Ecuaciones Diferenciales & 4 \\ \hline
MA-0793 & Tópicos de Geometría & 4 \\ \hline
MA-0794 & Tópicos de Modelación & 4 \\ \hline
MA-0795 & Tópicos de Computación Científica & 4 \\ \hline
MA-0796 & Tópicos de Álgebra & 4 \\ \hline
MA-0797 & Tópicos de Matemática Discreta & 4 \\ \hline
MA-0798 & Tópicos de Análisis de Datos & 4 \\ \hline
MA-0836 & Tópicos de Aprendizaje Automático & 4 \\ \hline
MA-0835 & Topología Algebraica & 4 \\ \hline
MA-0834 & Análisis Funcional & 4 \\ \hline
MA-0863 & Teoría Analítica de Números & 4 \\ \hline
MA-0862 & Análisis Armónico & 4 \\ \hline
MA-0861 & Teoría de Modelos & 4 \\ \hline
MA-0830 & Tópicos de Teoría de Números & 4 \\ \hline
MA-0831 & Tópicos de Lógica & 4 \\ \hline
MA-0832 & Tópicos de Topología & 4 \\ \hline
MA-0833 & Tópicos de Estadística & 4 \\ \hline
MA-0841 & Probabilidad & 4 \\ \hline
MA-0849 & Teoría de Módulos & 5 \\ \hline
MA-0858 & Teoría de Conjuntos & 4 \\ \hline
MA-0876 & Teoría Algebraica de Números & 4 \\ \hline
MA-0889 & Álgebra Conmutativa & 4 \\ \hline
MA-0920 & Ecuaciones en Derivadas Parciales Numéricas & 4 \\ \hline
MA-0928 & Procesos Estocásticos & 4 \\ \hline
\end{longtable}
\endgroup

{\footnotesize Los créditos de los cursos de Ciencias Actuariales (CA) se rigen por el plan de estudios de la carrera que los administra (---: no registrado en esta propuesta).\par}

\subsection*{Cursos optativos de otra disciplina (OPT-)}

\begingroup
\small
\arrayrulecolor{Azul-UCR!80}%
\renewcommand{\arraystretch}{1.2}%
\rowcolors{2}{white}{Celeste-UCR!12}%
\begin{longtable}{|l|>{\raggedright\arraybackslash}p{6cm}|c|>{\raggedright\arraybackslash}p{4.2cm}|}
\caption{Cursos optativos de otra disciplina}\label{res:opt-otra}\\
\hline
\rowcolor{Azul-UCR}
\textcolor{white}{\textbf{Sigla}} & \textcolor{white}{\textbf{Nombre del curso}} & \textcolor{white}{\textbf{Créditos}} & \textcolor{white}{\textbf{Unidad académica}} \\ \hline
\endfirsthead
\hline
\rowcolor{Azul-UCR}
\textcolor{white}{\textbf{Sigla}} & \textcolor{white}{\textbf{Nombre del curso}} & \textcolor{white}{\textbf{Créditos}} & \textcolor{white}{\textbf{Unidad académica}} \\ \hline
\endhead
\hline
\endfoot
B0106 & Biología General & 3 & Biología \\ \hline
B0107 & Laboratorio de Biología General & 1 & Biología \\ \hline
CI-0111 & Estructuras Discretas & 4 & Ciencias de la Computación e Informática \\ \hline
CI-0112 & Programación I & 4 & Ciencias de la Computación e Informática \\ \hline
CI-0113 & Programación II & 4 & Ciencias de la Computación e Informática \\ \hline
CI-0116 & Análisis de Algoritmos y Estructuras de Datos & 4 & Ciencias de la Computación e Informática \\ \hline
CI-0202 & Principios de Informática & 4 & Ciencias de la Computación e Informática \\ \hline
EC-1100 & Introducción a la Economía & 4 & Economía \\ \hline
EC-2100 & Teoría Microeconómica I & 4 & Economía \\ \hline
EC-3200 & Teoría Macroeconómica I & 4 & Economía \\ \hline
F0047 & Filosofía de la Física & 3 & Filosofía \\ \hline
F2001 & Seminario: Filosofía de la Física & 3 & Filosofía \\ \hline
F2002 & Filosofía de la Matemática & 3 & Filosofía \\ \hline
F2012 & Seminario de Epistemología & 3 & Filosofía \\ \hline
F2014 & Teoría del Método & 3 & Filosofía \\ \hline
F2019 & Seminario: Teorías y Revoluciones Científicas & 3 & Filosofía \\ \hline
F2022 & Filosofía de la Mente & 3 & Filosofía \\ \hline
F2334 & Filosofía de la Ciencia & 3 & Filosofía \\ \hline
F2454 & Introducción a la Epistemología & 3 & Filosofía \\ \hline
F2604 & Historia de la Ciencia & 3 & Filosofía \\ \hline
F4754 & Filosofía del Lenguaje & 3 & Filosofía \\ \hline
F5504 & Introducción a la Ética & 3 & Filosofía \\ \hline
FS-0210 & Física General I & 3 & Física \\ \hline
FS-0211 & Laboratorio de Física General I & 1 & Física \\ \hline
FS-0230 & Física I & 3 & Física \\ \hline
FS-0310 & Física General II & 3 & Física \\ \hline
FS-0311 & Laboratorio de Física General II & 1 & Física \\ \hline
FS-0330 & Física II & 3 & Física \\ \hline
FS-0331 & Laboratorio de Física Experimental I & 2 & Física \\ \hline
FS-0410 & Física General III & 3 & Física \\ \hline
FS-0411 & Laboratorio de Física General III & 1 & Física \\ \hline
FS-4101 & Introducción a la Meteorología & 3 & Física \\ \hline
G0111 & Fundamentos de Geología & 4 & Centroamericana de Geología \\ \hline
G0214 & Geología General & 4 & Centroamericana de Geología \\ \hline
QU-0100 & Química General I & 3 & Química \\ \hline
QU-0101 & Laboratorio de Química General I & 1 & Química \\ \hline
QU-0102 & Química General II & 3 & Química \\ \hline
QU-0103 & Laboratorio de Química General II & 1 & Química \\ \hline
QU-0200 & Química Analítica Cuantitativa I & 3 & Química \\ \hline
QU-0201 & Laboratorio de Química Analítica Cuantitativa I & 2 & Química \\ \hline
QU-0210 & Fundamentos de Química Orgánica & 4 & Química \\ \hline
QU-0211 & Laboratorio de Fundamentos de Química Orgánica & 1 & Química \\ \hline
XS-3150 & Diseño de Experimentos & 4 & Estadística \\ \hline
\end{longtable}
\endgroup


%%%%%%%%%%%%%%%%%%%%%%%%%%%%%%%%%%%%%%%%%%%%%%%%%%
\section{Descripción de cursos}

Para cada curso propio de la carrera se presenta su objetivo general y una síntesis de la temática. Los programas completos (objetivos específicos, contenidos, metodología, evaluación y bibliografía) se incluyen en la propuesta curricular. Los cursos de Ciencias Actuariales (CA) y los optativos de otra disciplina conservan los programas de las unidades que los administran.

% Generado por scripts/generate_resumen_ejecutivo.py. No editar a mano.

\begingroup
\small
\arrayrulecolor{Azul-UCR!80}%
\renewcommand{\arraystretch}{1.2}%
\rowcolors{2}{white}{Celeste-UCR!12}%
\begin{longtable}{|>{\raggedright\arraybackslash}p{3.3cm}|>{\raggedright\arraybackslash}p{5.5cm}|>{\raggedright\arraybackslash}p{6.9cm}|}
\caption{Descripción de los cursos propios de la carrera}\label{res:descripciones}\\
\hline
\rowcolor{Azul-UCR}
\textcolor{white}{\textbf{Curso (créditos)}} & \textcolor{white}{\textbf{Objetivo general}} & \textcolor{white}{\textbf{Temática resumida}} \\ \hline
\endfirsthead
\hline
\rowcolor{Azul-UCR}
\textcolor{white}{\textbf{Curso (créditos)}} & \textcolor{white}{\textbf{Objetivo general}} & \textcolor{white}{\textbf{Temática resumida}} \\ \hline
\endhead
\hline
\endfoot
\textbf{MA-0151}\newline Fundamentos de Álgebra, Trigonometría y Geometría Analítica (3) & Aplicar las nociones, intuiciones y conceptos básicos del álgebra, trigonometría y geometría analítica para la resolución de problemas. & Números reales y sus propiedades básicas; Trigonometría; Números complejos; Geometría vectorial y analítica básica; Proyectos de descubrimiento. \\ \hline
\textbf{MA-0152}\newline Matemática Exploratoria (3) & Iniciar al estudiantado en el quehacer matemático mediante la resolución de problemas de teoría de números, combinatoria y geometría, a través de la exploración, la conjetura y la justificación de resultados. & Introducción a los modos de trabajo y desarrollo de la matemática; Estrategias generales de resolución de problemas; Resolución de problemas concretos; Introducción a la formalización de cuantificadores. \\ \hline
\textbf{MA-0251}\newline Cálculo en una Variable (4) & Aplicar las nociones, intuiciones y conceptos básicos del cálculo diferencial e integral en una variable para la resolución de problemas. & Introducción a las funciones; Límites y continuidad; Derivación; Integración de Riemann. \\ \hline
\textbf{MA-0252}\newline Introducción a las Demostraciones (3) & Desarrollar en el estudiantado las habilidades de comprensión y aplicación de estrategias y métodos de demostración en matemáticas. & Lógica básica y estrategias de demostración; Teoría de conjuntos; Relaciones y funciones; Inducción matemática y recurrencia; Teoría de números. \\ \hline
\textbf{MA-0341}\newline Matemática Computacional (4) & Implementar el pensamiento algorítmico en la resolución de problemas matemáticos, mediante la ejecución de programas en lenguajes de programación de alto nivel. & Introducción a lenguajes de programación; Introducción a algoritmos; Visualizaciones; Algoritmos básicos en teoría de números; Algoritmos básicos en cálculo diferencial e integral; Algoritmos básicos en álgebra lineal y aplicaciones. \\ \hline
\textbf{MA-0351}\newline Cálculo en Varias Variables (4) & Aplicar las nociones, intuiciones y conceptos básicos del cálculo diferencial e integral en varias variables para la resolución de problemas. & Sucesiones y series; Derivación de funciones de varias variables; Integración en varias variables operacional. \\ \hline
\textbf{MA-0361}\newline Álgebra Lineal I (4) & Aplicar las nociones, intuiciones y conceptos básicos del álgebra lineal para la resolución de problemas. & Sistemas de ecuaciones lineales; Vectores y matrices; Determinantes; Espacios vectoriales; Transformaciones lineales. \\ \hline
\textbf{MA-0451}\newline Principios de Análisis I (4) & Aplicar los conceptos básicos de completitud, convergencia y diferenciación para la demostración de propiedades del análisis en una variable real. & Completitud y propiedades de los números reales; Sucesiones y series numéricas; Límites y continuidad; Diferenciación; Sucesiones y series de funciones; Construcción de las funciones trascendentes elementales. \\ \hline
\textbf{MA-0461}\newline Álgebra Lineal II (4) & Aplicar las nociones, intuiciones y conceptos del álgebra lineal para la resolución de problemas. & Polinomios; Dualidad; Valores y vectores propios; Espacios con producto interno; Triangularización, diagonalización y formas canónicas; Formas bilineales y multilineales. \\ \hline
\textbf{MA-0471}\newline Introducción a la Geometría Diferencial (3) & Analizar de forma rigurosa e intuitiva la geometría diferencial de curvas y superficies, mediante cálculos y demostraciones sobre sus propiedades locales y globales, tanto extrínsecas como intrínsecas. & Curvas; Superficies; Aplicación de Gauss; Geometría intrínseca de superficies; Geometría proyectiva. \\ \hline
\textbf{MA-0496}\newline Teoría de Números (3) & Aplicar los conceptos básicos de la teoría de números elemental para la resolución de problemas. & Repaso de divisibilidad en los enteros; Números primos; Congruencias; Funciones aritméticas; Raíces primitivas; Congruencias cuadráticas y la Ley de Reciprocidad cuadrática; Enteros en cuerpos de números cuadráticos. \\ \hline
\textbf{MA-0501}\newline Análisis Numérico I (4) & Aplicar la teoría de análisis numérico y sus herramientas básicas en la resolución de problemas y modelos matemáticos, mediante el estudio teórico de algoritmos y la implementación de algoritmos en un lenguaje de alto nivel. & Error numérico; Aproximación de funciones; Ecuaciones no lineales; Integración; Álgebra lineal numérica; Ecuaciones diferenciales ordinarias; Tema opcional. \\ \hline
\textbf{MA-0556}\newline Ecuaciones Diferenciales (4) & Establecer criterios para determinar la existencia, unicidad y estabilidad de soluciones de ecuaciones diferenciales ordinarias. & Introducción y ecuaciones diferenciales ordinarias de primer orden; Sistemas de ecuaciones diferenciales lineales; La transformada de Laplace; Teoría de existencia y unicidad; Continuación de soluciones; Dependencia de valores iniciales y parámetros; Tema a elegir por la persona docente. \\ \hline
\textbf{MA-0555}\newline Principios de Análisis II (4) & Aplicar los conceptos básicos de convergencia, diferenciación e integración para la demostración de propiedades del análisis en una y varias variables reales. & Integración de Riemann en $\mathbb{R}$; Topología métrica de $\mathbb{R}^n$; Funciones de varias variables; Diferenciación en varias variables; Cálculo y análisis vectorial. \\ \hline
\textbf{MA-0635}\newline Topología General (4) & Comprender los conceptos fundamentales de los espacios métricos y de los espacios topológicos, así como las funciones continuas entre ellos, y desarrollar herramientas, tanto conjuntistas como algebraicas, para analizar y distinguir dichos espacios. & Repaso de teoría de conjuntos; Espacios métricos; Espacios topológicos; Axiomas de separación; Homotopía y el grupo fundamental. \\ \hline
\textbf{MA-0571}\newline Álgebra Abstracta I (4) & Introducir a la persona estudiante a los conceptos y métodos propios del álgebra abstracta a través del estudio de la teoría de anillos y la teoría de módulos. & Definiciones, ejemplos y propiedades básicas de anillos; Morfismos, ideales y cocientes de anillos; Teoría de divisibilidad y factorización en anillos conmutativos; Anillos de polinomios y factorización; Definiciones, ejemplos y propiedades básicas de módulos; Grupos. \\ \hline
\textbf{MA-0661}\newline Álgebra Abstracta II (5) & Analizar y aplicar los conceptos fundamentales de la teoría de grupos, cuerpos y extensiones de cuerpos, para comprender la estructura de las extensiones algebraicas y su relación con los grupos de automorfismos, mediante el uso de resultados clásicos de álgebra. & Teoremas de estructura de grupos; Cuerpos; Automorfismos; Extensiones normales; Extensiones separables e inseparables; Teorema Fundamental de la Teoría de Galois; Cuerpos finitos. \\ \hline
\textbf{MA-0647}\newline Modelación Matemática (4) & Modelar problemas determinísticos a través de ecuaciones diferenciales y sistemas dinámicos para situaciones de ciencias, ingeniería o contextos interdisciplinarios, evaluando la validez del modelo y comunicando resultados con claridad. & El ciclo de modelación matemática; Modelación continua con ecuaciones diferenciales y sistemas dinámicos; Modelación discreta y simulación computacional; Validación, incertidumbre e interpretación con datos; Temas complementarios a elección de la persona docente; Proyecto final integrador (transversal, semanas 11--15). \\ \hline
\textbf{MA-0730}\newline Análisis Real I (5) & Aplicar los conceptos y resultados básicos del análisis real para la resolución de problemas. & La integral de Riemann-Stieltjes; Medida de Lebesgue en $\R^n$; Funciones Lebesgue medibles; La integral de Lebesgue; Espacios $L^p$. \\ \hline
\textbf{MA-0735}\newline Análisis Complejo (4) & Aplicar los conceptos y resultados básicos del análisis complejo para la resolución de problemas teóricos y prácticos. & Números complejos y funciones en el plano complejo; Diferenciación en el plano complejo; Integración compleja, el Teorema de Cauchy y sus aplicaciones; Singularidades, funciones meromorfas y el logaritmo complejo; Funciones enteras; Tópicos adicionales (cubrir al menos uno a criterio de la persona docente). \\ \hline
\textbf{MA-0880}\newline Comunicación en las Ciencias (2) & Que la persona estudiante sea capaz de comunicarse en forma oral y escrita para transmitir de forma efectiva ideas, resultados y argumentos matemáticos con claridad, rigor y pertinencia en contextos académicos y profesionales. & Comunicación científica y ética profesional; Escritura matemática y científica; Presentación oral y apoyos visuales; Búsqueda bibliográfica y gestión de fuentes; Comunicación en docencia y divulgación; Planificación de carrera y comunicación profesional; Presentaciones integradoras del estudiantado. \\ \hline
\textbf{MA-0881}\newline Práctica Profesional en Matemática Pura (6) & Que la persona estudiante integre y aplique conocimientos de matemática pura en un entorno académico o profesional, desarrollando autonomía, responsabilidad y capacidad de comunicar resultados. & Los contenidos no se fijan de manera permanente en el plan de estudios: dependen de la modalidad y del plan de trabajo aprobado para cada persona estudiante. \\ \hline
\textbf{MA-0883}\newline Práctica Profesional en Matemática Aplicada (6) & Que la persona estudiante integre y aplique conocimientos de matemática aplicada en un entorno académico o profesional, desarrollando autonomía, responsabilidad y capacidad de comunicar resultados. & Los contenidos no se fijan de manera permanente en el plan de estudios: dependen de la modalidad y del plan de trabajo aprobado para cada persona estudiante. \\ \hline
\textbf{MA-0982}\newline Seminario de Estudios Dirigidos en Matemática Pura I (7) & Que la persona estudiante desarrolle, de manera autónoma, actividades asociadas a la investigación en matemática pura. & Los contenidos no se fijan de manera permanente en el plan de estudios: dependen del tema y de los objetivos acordados para cada persona estudiante. \\ \hline
\textbf{MA-1081}\newline Seminario de Estudios Dirigidos en Matemática Pura II (7) & Que la persona estudiante desarrolle, de manera autónoma, actividades asociadas a la investigación en matemática pura. & Los contenidos no se fijan de manera permanente en el plan de estudios: dependen del tema y de los objetivos acordados para cada persona estudiante. \\ \hline
\textbf{MA-0984}\newline Seminario de Estudios Dirigidos en Matemática Aplicada I (7) & Que la persona estudiante desarrolle, de manera autónoma, actividades asociadas a la investigación en matemática aplicada. & Los contenidos no se fijan de manera permanente en el plan de estudios: dependen del tema y de los objetivos acordados para cada persona estudiante. \\ \hline
\textbf{MA-1082}\newline Seminario de Estudios Dirigidos en Matemática Aplicada II (7) & Que la persona estudiante desarrolle, de manera autónoma, actividades asociadas a la investigación en matemática aplicada. & Los contenidos no se fijan de manera permanente en el plan de estudios: dependen del tema y de los objetivos acordados para cada persona estudiante. \\ \hline
\textbf{MA-0526}\newline Combinatoria (4) & Desarrollar en la persona estudiante competencias para abordar problemas combinatorios mediante técnicas de conteo, teoría de grafos, métodos probabilísticos y combinatoria extremal. & Herramientas de conteo; Introducción a la teoría de grafos; Método probabilístico; Funciones generatrices; Combinatoria extremal. \\ \hline
\textbf{MA-0528}\newline Sistemas Dinámicos (4) & Que la persona estudiante adquiera las herramientas fundamentales para analizar sistemas dinámicos desde los puntos de vista topológico, diferenciable y ergódico, reconocer estructuras recurrentes, construir modelos simbólicos, estudiar medidas invariantes e interpretar el comportamiento estadístico de las órbitas mediante promedios temporales. & Sistemas dinámicos en tiempo discreto y continuo; Conjuntos límite, recurrencia y conjugaciones; Difeomorfismos del intervalo y linealización local; Codificación simbólica; Medidas invariantes y teoría ergódica; Mapas por ramas, mapa de Gauss y distorsión acotada; La familia cuadrática. \\ \hline
\textbf{MA-0838}\newline Geometría Diferencial (4) & Comprender y aplicar las herramientas fundamentales del cálculo diferencial e integral y de la geometría riemanniana en variedades, demostrando resultados centrales. & Variedades diferenciales y topología diferencial; Formas diferenciales; Integración en variedades; Geometría riemanniana. \\ \hline
\textbf{MA-0837}\newline Geometría Algebraica (4) & Introducir a la persona estudiante a los conceptos, métodos y resultados fundamentales de la Geometría Algebraica clásica, con énfasis en variedades y curvas algebraicas, sus singularidades, propiedades proyectivas y birracionales, y sus invariantes geométricos y aritméticos. \ & Conjuntos algebraicos afines y el Teorema de los Ceros de Hilbert; Variedades afines y propiedades locales; Propiedades locales de curvas planas; Geometría proyectiva; Curvas proyectivas planas e intersecciones; Morfismos y geometría birracional; Resolución de singularidades de curvas; Divisores, género y Teorema de Riemann--Roch; Curvas sobre cuerpos finitos, funciones zeta y conjeturas de Weil. \\ \hline
\textbf{MA-0734}\newline Optimización (4) & Desarrollar los fundamentos matemáticos de la optimización para establecer la existencia y caracterizar las soluciones de problemas con y sin restricciones. & Fundamentos y optimización sin restricciones; Programación lineal; Optimización con restricciones de igualdad y desigualdad; Proyectos de descubrimiento. \\ \hline
\textbf{MA-0756}\newline Ecuaciones Diferenciales Parciales (5) & Desarrollar competencias teóricas para formular, analizar y resolver problemas de ecuaciones diferenciales parciales mediante métodos clásicos y formulaciones débiles. & Ecuaciones clásicas fundamentales; Espacios de Sobolev; Ecuaciones elípticas de segundo orden; Ecuaciones de evolución lineales. \\ \hline
\textbf{MA-0781}\newline Lógica (5) & Desarrollar los conocimientos y herramientas fundamentales de la lógica matemática para analizar y formalizar razonamientos matemáticos y demostrar los principales resultados del cálculo de predicados de primer orden. & Ordinales y cardinales; Lógica de primer orden; Teorema de completitud y teorema de compacidad; Nociones básicas de teorías y estructuras. \\ \hline
\textbf{MA-0783}\newline Integración (4) & Desarrollar competencias en teoría abstracta de la medida e integración y en sus aplicaciones clásicas dentro del análisis matemático. & Repaso de integración de Lebesgue; Curvas rectificables e isoperimetría; Teoría abstracta de la medida e integración; Teoremas ergódicos; Medidas de Hausdorff; Medidas de Haar; Transporte óptimo. \\ \hline
\textbf{MA-0788}\newline Análisis Real II (4) & Aplicar los conceptos y resultados intermedios del Análisis Real para la resolución de problemas. & Diferenciación; Continuidad absoluta y el Teorema Fundamental del Cálculo; Espacios de Banach; Espacios de Hilbert; Introducción al Análisis de Fourier. \\ \hline
\textbf{MA-0835}\newline Topología Algebraica (4) & Que la persona estudiante domine los conceptos y técnicas elementales de la topología algebraica y pueda aplicarlos al análisis de espacios y aplicaciones continuas. & Homotopía (noción básica); Grupo fundamental; Homología; Cohomología de haces (Sheaf Cohomology); Homotopía superior. \\ \hline
\textbf{MA-0834}\newline Análisis Funcional (4) & Que la persona estudiante sea capaz de resolver problemas y desarrollar teoría a partir de los fundamentos del análisis funcional en espacios de Banach y de Hilbert. & Espacios de Banach y espacios $L^p$; Espacios de Hilbert; Operadores lineales acotados; Teoremas fundamentales; Dualidad y topologías débiles; Operadores compactos y teorema espectral. \\ \hline
\textbf{MA-0863}\newline Teoría Analítica de Números (4) & Introducir a la persona estudiante a las principales ideas, técnicas y resultados de la Teoría Analítica de Números, enfatizando el uso de herramientas del análisis real y complejo para estudiar funciones aritméticas y la distribución de los números primos. \ & Funciones aritméticas y convolución de Dirichlet; Funciones sumatorias y órdenes promedio; Series de Dirichlet y herramientas de análisis complejo; La función zeta de Riemann; Ceros de la función zeta y distribución de los números primos; Caracteres y funciones $L$ de Dirichlet; Primos en progresiones aritméticas; Problemas abiertos y desarrollos posteriores. \\ \hline
\textbf{MA-0862}\newline Análisis Armónico (4) & Que la persona estudiante sea capaz de resolver problemas y desarrollar teoría a partir de los fundamentos del análisis armónico de variable real en $\R^n$. & Espacios $L^p$; La transformada de Fourier; Funciones maximales; Las transformadas de Hilbert y de Riesz; Integrales singulares de tipo convolución; Teoría de Littlewood--Paley. \\ \hline
\textbf{MA-0861}\newline Teoría de Modelos (4) & Comprender los principios básicos de la teoría de modelos y de sus aplicaciones a las estructuras algebraicas. & Fundamentos de lógica de primer orden; Principios básicos de teoría de modelos; Eliminación de cuantificadores; Aplicaciones al álgebra y la geometría algebraica; Tipos y compacidad. \\ \hline
\textbf{MA-0841}\newline Probabilidad (4) & Desarrollar los fundamentos de la teoría de la probabilidad en el marco de la teoría de la medida para modelar incertidumbre, demostrar resultados centrales y aplicarlos en problemas matemáticos e interdisciplinarios. & Probabilidad elemental; Teoría de la medida e integración; Teoría de probabilidad; Independencia; Convergencia de variables aleatorias; Probabilidad condicional. \\ \hline
\textbf{MA-0849}\newline Teoría de Módulos (5) & Comprender los fundamentos de la teoría de módulos y sus principales aplicaciones algebraicas. & Módulos; Construcciones básicas; Módulos noetherianos y artinianos; Módulos sobre dominios de ideales principales; Módulos proyectivos, inyectivos y planos; Producto tensorial; Introducción al álgebra homológica (Optativo). \\ \hline
\textbf{MA-0858}\newline Teoría de Conjuntos (4) & Introducir al estudiantado en la teoría axiomática de conjuntos de Zermelo--Fraenkel, desarrollando las herramientas fundamentales como cardinales, ordinales y aritmética transfinita. & Motivación y antecedentes históricos; El sistema axiomático de Zermelo--Fraenkel; Relaciones y funciones; Órdenes; Los números naturales; Números cardinales; Números ordinales; El axioma de elección y la hipótesis del continuo. \\ \hline
\textbf{MA-0876}\newline Teoría Algebraica de Números (4) & Introducir a la persona estudiante a los conceptos, métodos y ejemplos fundamentales de la Teoría de Números Algebraicos, enfatizando la aritmética de los anillos de enteros de cuerpos de números, la factorización de ideales, el grupo de clases, el grupo de unidades, la factorización de primos en extensiones y el uso de herramientas computacionales para explorar ejemplos concretos. & Cuerpos de números y enteros algebraicos; Ideales y factorización; Factorización de ideales primos y ramificación; El grupo de clases y el teorema de Minkowski; Unidades y teorema de Dirichlet; Conteo de ideales y función zeta de Dedekind; Teoría de descomposición e inercia en extensiones de Galois. \\ \hline
\textbf{MA-0889}\newline Álgebra Conmutativa (4) & Introducir a la persona estudiante a los conceptos, métodos y resultados fundamentales del Álgebra Conmutativa, enfatizando las propiedades de anillos y módulos noetherianos, la localización, la dependencia integral, la teoría de dimensión, los métodos computacionales y la estructura local de los anillos conmutativos. & Módulos y herramientas algebraicas fundamentales; El diccionario entre álgebra y geometría; Noetherianidad y descomposición primaria; Localización y anillos locales; Dependencia integral y normalización; Dimensión de Krull; Bases de Gröbner y métodos computacionales; Anillos locales regulares y anillos de dimensión uno. \\ \hline
\textbf{MA-0920}\newline Ecuaciones en Derivadas Parciales Numéricas (4) & Aplicar la teoría del análisis numérico y sus herramientas básicas para aproximar soluciones numéricas de ecuaciones diferenciales parciales. & Introducción a las ecuaciones diferenciales parciales; Método de diferencias finitas; Método de elemento finito; Resolución numérica de sistemas lineales; Resolución numérica de la ecuación de calor y de onda; Temas adicionales. \\ \hline
\textbf{MA-0477}\newline Ecuaciones Diferenciales Parciales Aplicadas (4) & Aplicar métodos clásicos del análisis de EDP para resolver y analizar ecuaciones lineales que surgen en modelos matemáticos. & EDP de primer orden; Ecuaciones de Laplace y Poisson; Ecuación del calor; Ecuación de onda; Introducción al cálculo de variaciones. \\ \hline
\textbf{MA-0759}\newline Tópicos de Matemática Financiera (4) & Profundizar en un tema de la matemática financiera, definido para cada oferta del curso, mediante el estudio guiado de resultados, problemas y, cuando corresponda, aplicaciones. & Los contenidos no se fijan de manera permanente en el plan de estudios: son abiertos por oferta, siempre que pertenezcan al ámbito de Matemática financiera. Al inicio del ciclo lectivo, la persona docente debe comunicar por escrito al estudiantado. \\ \hline
\textbf{MA-0790}\newline Tópicos de Análisis (4) & Profundizar en un tema de análisis matemático, definido para cada oferta del curso, mediante el estudio guiado de resultados, problemas y, cuando corresponda, aplicaciones. & Los contenidos no se fijan de manera permanente en el plan de estudios: son abiertos por oferta, siempre que pertenezcan al ámbito del Análisis matemático. \\ \hline
\textbf{MA-0791}\newline Tópicos de Probabilidad (4) & Profundizar en un tema de la probabilidad, definido para cada oferta del curso, mediante el estudio guiado de resultados, problemas y, cuando corresponda, aplicaciones. & Los contenidos no se fijan de manera permanente en el plan de estudios: son abiertos por oferta, siempre que pertenezcan al ámbito de Probabilidad. Al inicio del ciclo lectivo, la persona docente debe comunicar por escrito al estudiantado. \\ \hline
\textbf{MA-0792}\newline Tópicos de Ecuaciones Diferenciales (4) & Profundizar en un tema de las ecuaciones diferenciales, definido para cada oferta del curso, mediante el estudio guiado de resultados, problemas y, cuando corresponda, aplicaciones. & Los contenidos no se fijan de manera permanente en el plan de estudios: son abiertos por oferta, siempre que pertenezcan al ámbito de Ecuaciones diferenciales. Al inicio del ciclo lectivo, la persona docente debe comunicar por escrito al estudiantado. \\ \hline
\textbf{MA-0793}\newline Tópicos de Geometría (4) & Profundizar en un tema de la geometría, definido para cada oferta del curso, mediante el estudio guiado de resultados, problemas y, cuando corresponda, aplicaciones. & Los contenidos no se fijan de manera permanente en el plan de estudios: son abiertos por oferta, siempre que pertenezcan al ámbito de Geometría. Al inicio del ciclo lectivo, la persona docente debe comunicar por escrito al estudiantado. \\ \hline
\textbf{MA-0794}\newline Tópicos de Modelación (4) & Profundizar en un tema de la modelación matemática, definido para cada oferta del curso, mediante el estudio guiado de resultados, problemas y, cuando corresponda, aplicaciones. & Los contenidos no se fijan de manera permanente en el plan de estudios: son abiertos por oferta, siempre que pertenezcan al ámbito de Modelación matemática. Al inicio del ciclo lectivo, la persona docente debe comunicar por escrito al estudiantado. \\ \hline
\textbf{MA-0795}\newline Tópicos de Computación Científica (4) & Profundizar en un tema de la computación científica, definido para cada oferta del curso, mediante el estudio guiado de resultados, problemas y, cuando corresponda, aplicaciones. & Los contenidos no se fijan de manera permanente en el plan de estudios: son abiertos por oferta, siempre que pertenezcan al ámbito de Computación científica. Al inicio del ciclo lectivo, la persona docente debe comunicar por escrito al estudiantado. \\ \hline
\textbf{MA-0796}\newline Tópicos de Álgebra (4) & Profundizar en un tema del álgebra, definido para cada oferta del curso, mediante el estudio guiado de resultados, problemas y, cuando corresponda, aplicaciones. & Los contenidos no se fijan de manera permanente en el plan de estudios: son abiertos por oferta, siempre que pertenezcan al ámbito de Álgebra. Al inicio del ciclo lectivo, la persona docente debe comunicar por escrito al estudiantado. \\ \hline
\textbf{MA-0797}\newline Tópicos de Matemática Discreta (4) & Profundizar en un tema de la matemática discreta, definido para cada oferta del curso, mediante el estudio guiado de resultados, problemas y, cuando corresponda, aplicaciones. & Los contenidos no se fijan de manera permanente en el plan de estudios: son abiertos por oferta, siempre que pertenezcan al ámbito de Matemática discreta. Al inicio del ciclo lectivo, la persona docente debe comunicar por escrito al estudiantado. \\ \hline
\textbf{MA-0798}\newline Tópicos de Análisis de Datos (4) & Profundizar en un tema del análisis de datos, definido para cada oferta del curso, mediante el estudio guiado de resultados, problemas y, cuando corresponda, aplicaciones. & Los contenidos no se fijan de manera permanente en el plan de estudios: son abiertos por oferta, siempre que pertenezcan al ámbito de Análisis de datos. Al inicio del ciclo lectivo, la persona docente debe comunicar por escrito al estudiantado. \\ \hline
\textbf{MA-0836}\newline Tópicos de Aprendizaje Automático (4) & Profundizar en un tema del aprendizaje automático, definido para cada oferta del curso, mediante el estudio guiado de resultados, problemas y, cuando corresponda, aplicaciones. & Los contenidos no se fijan de manera permanente en el plan de estudios: son abiertos por oferta, siempre que pertenezcan al ámbito de Aprendizaje automático. Al inicio del ciclo lectivo, la persona docente debe comunicar por escrito al estudiantado. \\ \hline
\textbf{MA-0830}\newline Tópicos de Teoría de Números (4) & Profundizar en un tema de la teoría de números, definido para cada oferta del curso, mediante el estudio guiado de resultados, problemas y, cuando corresponda, aplicaciones. & Los contenidos no se fijan de manera permanente en el plan de estudios: son abiertos por oferta, siempre que pertenezcan al ámbito de Teoría de números. Al inicio del ciclo lectivo, la persona docente debe comunicar por escrito al estudiantado. \\ \hline
\textbf{MA-0831}\newline Tópicos de Lógica (4) & Profundizar en un tema de la lógica, definido para cada oferta del curso, mediante el estudio guiado de resultados, problemas y, cuando corresponda, aplicaciones. & Los contenidos no se fijan de manera permanente en el plan de estudios: son abiertos por oferta, siempre que pertenezcan al ámbito de Lógica. Al inicio del ciclo lectivo, la persona docente debe comunicar por escrito al estudiantado. \\ \hline
\textbf{MA-0832}\newline Tópicos de Topología (4) & Profundizar en un tema de la topología, definido para cada oferta del curso, mediante el estudio guiado de resultados, problemas y, cuando corresponda, aplicaciones. & Los contenidos no se fijan de manera permanente en el plan de estudios: son abiertos por oferta, siempre que pertenezcan al ámbito de Topología. Al inicio del ciclo lectivo, la persona docente debe comunicar por escrito al estudiantado. \\ \hline
\textbf{MA-0833}\newline Tópicos de Estadística (4) & Profundizar en un tema de la estadística, definido para cada oferta del curso, mediante el estudio guiado de resultados, problemas y, cuando corresponda, aplicaciones. & Los contenidos no se fijan de manera permanente en el plan de estudios: son abiertos por oferta, siempre que pertenezcan al ámbito de Estadística. Al inicio del ciclo lectivo, la persona docente debe comunicar por escrito al estudiantado. \\ \hline
\textbf{MA-0928}\newline Procesos Estocásticos (4) & Manejar los conceptos y técnicas básicas de la teoría de procesos estocásticos en tiempo continuo, con especial atención al movimiento browniano y sus aplicaciones. & Procesos gaussianos; Movimiento browniano; Filtraciones, martingalas y semimartingalas continuas; Variación cuadrática e integración estocástica; Fórmula de Itô y consecuencias; Procesos de Markov; Movimiento browniano y ecuaciones en derivadas parciales; Ecuaciones diferenciales estocásticas. \\ \hline
\end{longtable}
\endgroup


%%%%%%%%%%%%%%%%%%%%%%%%%%%%%%%%%%%%%%%%%%%%%%%%%%
\section{Tabla de correspondencia del equipo docente}

La tabla presenta, para cada curso, el personal docente del Departamento de Matemática disponible para impartirlo y su último grado académico. Se espera que la gran mayoría de los cursos sean impartidos por docentes con doctorado; se considerará la asignación de docentes con maestría para cursos de índole general o afines a su área de trabajo. Un mismo docente puede figurar en varios cursos.

% Generado por scripts/generate_resumen_ejecutivo.py. No editar a mano.

\begingroup
\small
\arrayrulecolor{Azul-UCR!80}%
\renewcommand{\arraystretch}{1.2}%
\begin{longtable}{|>{\raggedright\arraybackslash}p{6.9cm}|>{\raggedright\arraybackslash}p{9.8cm}|}
\caption{Listado de las actividades de formación (cursos) y personal docente a cargo}\label{res:docentes}\\
\hline
\rowcolor{Azul-UCR}
\textcolor{white}{\textbf{Nombre del curso}} & \textcolor{white}{\textbf{Persona(s) docente(s) a cargo}} \\ \hline
\endfirsthead
\hline
\rowcolor{Azul-UCR}
\textcolor{white}{\textbf{Nombre del curso}} & \textcolor{white}{\textbf{Persona(s) docente(s) a cargo}} \\ \hline
\endhead
\hline
\endfoot
\rowcolor{Celeste-UCR!35}\multicolumn{2}{|l|}{\textbf{Escuela de Matemática, Departamento de Matemática}} \\ \hline
\textbf{MA-0151} Fundamentos de Álgebra, Trigonometría y Geometría Analítica & Villalobos Arias Mario; Guier Acosta Jorge; Sánchez Guevara Jesús; Acuña Larios Jennifer; Lacy Mora Allan \\ \hline
\textbf{MA-0152} Matemática Exploratoria & Calvo Alpízar Juan Gabriel; Zamora Luna Oscar; Mora Sáenz Germán José; Ulloa Esquivel Roberto; Bustamante Medina Ronald \\ \hline
\textbf{MA-0251} Cálculo en una Variable & Acuña Larios Jennifer; Campos Fernández José David; Ramírez González José Alexander; Zamora Luna Oscar; Alfaro Murillo Jorge; Fonseca Mora Christian \\ \hline
\textbf{MA-0252} Introducción a las Demostraciones & Montenegro Guzmán Samaria; Zamora Calero Rafael; Loría Sorio Jennifer de los Ángeles; Ulloa Esquivel Roberto; Mora Sáenz Germán José \\ \hline
\textbf{MA-0341} Matemática Computacional & Calvo Alpízar Juan Gabriel; Rodríguez Rojas Oldemar; Segura Ugalde Esteban; Alfaro Murillo Jorge; Guevara Villalobos Álvaro \\ \hline
\textbf{MA-0351} Cálculo en Varias Variables & Acuña Larios Jennifer; Bolaños Guerrero Raúl Esteban; Jiménez López David; Mena Arias Darío Alberto; Chaves Aguilar Santiago \\ \hline
\textbf{MA-0361} Álgebra Lineal I & Barquero Sánchez Adrián; Bustamante Medina Ronald; Lacy Mora Allan; Rosales Ortega José; Jiménez López David \\ \hline
\textbf{MA-0451} Principios de Análisis I & Acuña Valverde Luis; Cambronero Villalobos Santiago; Campos Fernández José David; Sánchez Chavarría Adriana; Bolaños Guerrero Raúl Esteban \\ \hline
\textbf{MA-0461} Álgebra Lineal II & Guier Acosta Jorge; Hernández Alvarado Alberto; Segura Ugalde Esteban; Campos Salas Daniel; Mora Sáenz Germán José \\ \hline
\textbf{MA-0471} Introducción a la Geometría Diferencial & Rosales Ortega José; Zúñiga Rojas Ronald Alberto; Ugalde Gómez William Javier; Chaves Aguilar Santiago \\ \hline
\textbf{MA-0496} Teoría de Números & Lacy Mora Allan; Coto Mora Edward; Barquero Sánchez Adrián; Ávila Artavia Josué David \\ \hline
\textbf{MA-0556} Ecuaciones Diferenciales & Montalto Cruz Carlos; Loría Sorio Jennifer de los Ángeles; Sánchez Chavarría Adriana; Campos Salas Daniel \\ \hline
\textbf{MA-0555} Principios de Análisis II & Bolaños Guerrero Raúl Esteban; Ugalde Gómez William Javier; Villalobos Alvarado Jorge; Cambronero Villalobos Santiago \\ \hline
\textbf{MA-0635} Topología General & Zamora Calero Rafael; Tall El Hadji Yaya; Villalobos Alvarado Jorge; Rosales Ortega José \\ \hline
\textbf{MA-0571} Álgebra Abstracta I & Barquero Sánchez Adrián; Hernández Alvarado Alberto; Bustamante Medina Ronald; Zúñiga Rojas Ronald Alberto \\ \hline
\textbf{MA-0661} Álgebra Abstracta II & Guier Acosta Jorge; Montenegro Guzmán Samaria; Ávila Artavia Josué David; Coto Mora Edward \\ \hline
\textbf{MA-0647} Modelación Matemática & Montalto Cruz Carlos; Sánchez Peña Fabio; Alfaro Murillo Jorge; Loría Sorio Jennifer de los Ángeles \\ \hline
\textbf{MA-0735} Análisis Complejo & Bolaños Guerrero Raúl Esteban; Méndez Hernández Pedro; Tall El Hadji Yaya; Gutiérrez Pavón Jonathan; Calvo Alpízar Juan Gabriel \\ \hline
\textbf{MA-0730} Análisis Real I & Acuña Valverde Luis; Mena Arias Darío Alberto; Mora Sáenz Germán José; Campos Fernández José David \\ \hline
\textbf{MA-0880} Comunicación en las Ciencias & Mena Arias Darío Alberto; Sánchez Chavarría Adriana; Sánchez Peña Fabio; Alfaro Murillo Jorge \\ \hline
\textbf{MA-0788} Análisis Real II & Cambronero Villalobos Santiago; Ramírez González José Alexander; Fonseca Mora Christian; Méndez Hernández Pedro \\ \hline
\textbf{MA-0876} Teoría Algebraica de Números & Coto Mora Edward; Lacy Mora Allan; Chaves Aguilar Santiago \\ \hline
\textbf{MA-0863} Teoría Analítica de Números & Barquero Sánchez Adrián; Ávila Artavia Josué David; Ulloa Esquivel Roberto \\ \hline
\textbf{MA-0858} Teoría de Conjuntos & Zamora Calero Rafael; Sánchez Guevara Jesús \\ \hline
\textbf{MA-0526} Combinatoria & Zamora Luna Oscar; Tall El Hadji Yaya \\ \hline
\textbf{MA-0528} Sistemas Dinámicos & Sánchez Chavarría Adriana; Tall El Hadji Yaya \\ \hline
\textbf{MA-0783} Integración & Gutiérrez Pavón Jonathan; Mena Arias Darío Alberto; Ugalde Gómez William Javier \\ \hline
\textbf{MA-0838} Geometría Diferencial & Ramírez González José Alexander; Rosales Ortega José; Ugalde Gómez William Javier \\ \hline
\textbf{MA-0837} Geometría Algebraica & Zúñiga Rojas Ronald Alberto; Coto Mora Edward; Ulloa Esquivel Roberto \\ \hline
\textbf{MA-0781} Lógica & Montenegro Guzmán Samaria; Zamora Calero Rafael; Guier Acosta Jorge \\ \hline
\textbf{MA-0734} Optimización & Acuña Larios Jennifer; Guevara Villalobos Álvaro; Segura Ugalde Esteban \\ \hline
\textbf{MA-0755} Ecuaciones Diferenciales Parciales & Méndez Hernández Pedro; Montalto Cruz Carlos; Acuña Valverde Luis; Campos Salas Daniel \\ \hline
\textbf{MA-0835} Topología Algebraica & Zúñiga Rojas Ronald Alberto; Ugalde Gómez William Javier; Ávila Artavia Josué David \\ \hline
\textbf{MA-0834} Análisis Funcional & Cambronero Villalobos Santiago; Jiménez López David; Villalobos Alvarado Jorge \\ \hline
\textbf{MA-0862} Análisis Armónico & Jiménez López David; Mena Arias Darío Alberto; Tall El Hadji Yaya \\ \hline
\textbf{MA-0861} Teoría de Modelos & Bustamante Medina Ronald; Guier Acosta Jorge; Montenegro Guzmán Samaria \\ \hline
\textbf{MA-0841} Probabilidad & Acuña Valverde Luis; Campos Fernández José David; Fonseca Mora Christian \\ \hline
\textbf{MA-0849} Teoría de Módulos & Hernández Alvarado Alberto; Lacy Mora Allan; Sánchez Guevara Jesús \\ \hline
\textbf{MA-0889} Álgebra Conmutativa & Hernández Alvarado Alberto; Ávila Artavia Josué David; Ulloa Esquivel Roberto \\ \hline
\textbf{MA-0920} Ecuaciones en Derivadas Parciales Numéricas & Calvo Alpízar Juan Gabriel; Segura Ugalde Esteban; Jiménez López David; Villalobos Arias Mario \\ \hline
\textbf{MA-0477} Ecuaciones Diferenciales Parciales Aplicadas & Segura Ugalde Esteban; Mora Sáenz Germán José; Villalobos Arias Mario \\ \hline
\textbf{MA-0790} Tópicos de Análisis & Acuña Valverde Luis; Cambronero Villalobos Santiago; Campos Salas Daniel \\ \hline
\textbf{MA-0791} Tópicos de Probabilidad & Acuña Larios Jennifer; Bolaños Guerrero Raúl Esteban; Gutiérrez Pavón Jonathan; Fonseca Mora Christian \\ \hline
\textbf{MA-0792} Tópicos de Ecuaciones Diferenciales & Méndez Hernández Pedro; Sánchez Peña Fabio; Ramírez González José Alexander \\ \hline
\textbf{MA-0793} Tópicos de Geometría & Rosales Ortega José; Chaves Aguilar Santiago; Villalobos Alvarado Jorge \\ \hline
\textbf{MA-0794} Tópicos de Modelación & Montalto Cruz Carlos; Sánchez Peña Fabio; Alfaro Murillo Jorge; Loría Sorio Jennifer de los Ángeles \\ \hline
\textbf{MA-0796} Tópicos de Álgebra & Sánchez Guevara Jesús; Chaves Aguilar Santiago; Hernández Alvarado Alberto \\ \hline
\textbf{MA-0797} Tópicos de Matemática Discreta & Méndez Hernández Pedro; Zamora Luna Oscar; Campos Salas Daniel; Villalobos Arias Mario \\ \hline
\textbf{MA-0830} Tópicos de Teoría de Números & Barquero Sánchez Adrián; Zamora Luna Oscar; Coto Mora Edward \\ \hline
\textbf{MA-0831} Tópicos de Lógica & Bustamante Medina Ronald; Sánchez Guevara Jesús; Montenegro Guzmán Samaria; Zamora Calero Rafael \\ \hline
\textbf{MA-0832} Tópicos de Topología & Zúñiga Rojas Ronald Alberto; Villalobos Alvarado Jorge; Rodríguez Rojas Oldemar \\ \hline
\textbf{MA-0928} Procesos Estocásticos & Campos Fernández José David; Gutiérrez Pavón Jonathan; Fonseca Mora Christian \\ \hline
\textbf{MA-0795} Tópicos de Computación Científica & Calvo Alpízar Juan Gabriel; Rodríguez Rojas Oldemar; Guevara Villalobos Álvaro \\ \hline
\textbf{MA-0798} Tópicos de Análisis de Datos & Rodríguez Rojas Oldemar; Sánchez Chavarría Adriana; Sánchez Peña Fabio \\ \hline
\textbf{MA-0836} Tópicos de Aprendizaje Automático & Guevara Villalobos Álvaro; Rodríguez Rojas Oldemar; Loría Sorio Jennifer de los Ángeles \\ \hline
\textbf{MA-0759} Tópicos de Matemática Financiera & Guevara Villalobos Álvaro; Ramírez González José Alexander; Villalobos Arias Mario \\ \hline
\rowcolor{Celeste-UCR!35}\multicolumn{2}{|l|}{\textbf{Escuela de Matemática, Departamento de Ciencias Actuariales}} \\ \hline
\textbf{CA-0204} Herramientas de Ciencia de Datos I & \textit{Curso ya existente, a cargo del personal docente que actualmente lo imparte.} \\ \hline
\textbf{CA-0721} Probabilidad & \textit{Curso ya existente, a cargo del personal docente que actualmente lo imparte.} \\ \hline
\textbf{CA-0303} Estadística Actuarial I & \textit{Curso ya existente, a cargo del personal docente que actualmente lo imparte.} \\ \hline
\textbf{CA-0304} Herramientas de Ciencia de Datos II & \textit{Curso ya existente, a cargo del personal docente que actualmente lo imparte.} \\ \hline
\textbf{CA-0411a} Análisis de Datos II & \textit{Curso ya existente, a cargo del personal docente que actualmente lo imparte.} \\ \hline
\textbf{CA-0403} Estadística Actuarial II & \textit{Curso ya existente, a cargo del personal docente que actualmente lo imparte.} \\ \hline
\textbf{CA-0512} Análisis de Datos II & \textit{Curso ya existente, a cargo del personal docente que actualmente lo imparte.} \\ \hline
\textbf{CA-0612} Series de Tiempo & \textit{Curso ya existente, a cargo del personal docente que actualmente lo imparte.} \\ \hline
\rowcolor{Celeste-UCR!35}\multicolumn{2}{|l|}{\textbf{Escuela de Matemática, otros cursos ya existentes}} \\ \hline
\textbf{MA-0501} Análisis Numérico I & \textit{Curso ya existente, a cargo del personal docente que actualmente lo imparte.} \\ \hline
\textbf{MA-0602} Introducción a la Optimización & \textit{Curso ya existente, a cargo del personal docente que actualmente lo imparte.} \\ \hline
\end{longtable}
\endgroup



%%%%%%%%%%%%%%%%%%%%%%%%%%%%%%%%%%%%%%%%%%%%%%%%%%
\section{Otros aspectos}

\subsection*{Plan de transición}
El nuevo plan entra en vigencia en 2028. El plan 02 se mantiene activo hasta 2030 y se inactiva posteriormente; quienes no lo hayan concluido al año 2031 inclusive pasan de oficio al plan 2028, con las convalidaciones de la tabla correspondiente. La estrategia se define por cohorte (C5 y anteriores, C6, C7, y C8 en adelante, que ingresa directamente al nuevo plan), de acuerdo con la Circular VD-18-2015. La proyección de carga docente en transición es de entre 10 y 12 tiempos completos por año (2028--2031).

\subsection*{Recursos y factibilidad}
La Escuela cuenta con aulas y oficinas en el edificio de Física-Matemática y en el edificio del CIMPA, dos laboratorios de cómputo, una licencia de software por estudiante y acceso bajo demanda a recursos bibliográficos (SIBDI). El personal administrativo es el que ya atiende la carrera vigente. Las prácticas profesionales se ofrecen en modalidad académica o profesional, con apoyo de la Vicerrectoría de Vida Estudiantil para pasantías e intercambio estudiantil.

\subsection*{Seguimiento y evaluación}
El Departamento de Matemática estará a cargo de la implementación. Se propone una comisión de evaluación curricular que dé seguimiento desde el primer ciclo de vigencia (revisión anual de guías metodológicas, valoración de cursos al final de cada ciclo, análisis de la evaluación docente, consulta bianual sobre la carga académica, actualización anual de indicadores de admisión, permanencia, rezago y graduación, y autoevaluación). Se prevé un diagnóstico del perfil de egreso a más tardar cinco años después de la implementación completa, y una evaluación integral con asesoría del CEA seis años después de la implementación, con consulta a personas graduadas y empleadoras. 

\section{Currículos vitae de personas docentes}

\textcolor{red}{Pendiente}

\end{document}
